% Options for packages loaded elsewhere
\PassOptionsToPackage{unicode}{hyperref}
\PassOptionsToPackage{hyphens}{url}
%
\documentclass[
  11pt,
]{article}
\usepackage{amsmath,amssymb}
\usepackage{iftex}
\ifPDFTeX
  \usepackage[T1]{fontenc}
  \usepackage[utf8]{inputenc}
  \usepackage{textcomp} % provide euro and other symbols
\else % if luatex or xetex
  \usepackage{unicode-math} % this also loads fontspec
  \defaultfontfeatures{Scale=MatchLowercase}
  \defaultfontfeatures[\rmfamily]{Ligatures=TeX,Scale=1}
\fi
\usepackage{lmodern}
\ifPDFTeX\else
  % xetex/luatex font selection
  % Native editor uses its default monospaced font.
\fi
% Use upquote if available, for straight quotes in verbatim environments
\IfFileExists{upquote.sty}{\usepackage{upquote}}{}
\IfFileExists{microtype.sty}{% use microtype if available
  \usepackage[]{microtype}
  \UseMicrotypeSet[protrusion]{basicmath} % disable protrusion for tt fonts
}{}
\makeatletter
\@ifundefined{KOMAClassName}{% if non-KOMA class
  \IfFileExists{parskip.sty}{%
    \usepackage{parskip}
  }{% else
    \setlength{\parindent}{0pt}
    \setlength{\parskip}{6pt plus 2pt minus 1pt}}
}{% if KOMA class
  \KOMAoptions{parskip=half}}
\makeatother
\usepackage{xcolor}
\usepackage[margin=1in]{geometry}
\usepackage{color}
\usepackage{fancyvrb}
\newcommand{\VerbBar}{|}
\newcommand{\VERB}{\Verb[commandchars=\\\{\}]}
\DefineVerbatimEnvironment{Highlighting}{Verbatim}{commandchars=\\\{\}}
% Add ',fontsize=\small' for more characters per line
\newenvironment{Shaded}{}{}
\newcommand{\AlertTok}[1]{\textcolor[rgb]{1.00,0.00,0.00}{\textbf{#1}}}
\newcommand{\AnnotationTok}[1]{\textcolor[rgb]{0.38,0.63,0.69}{\textbf{\textit{#1}}}}
\newcommand{\AttributeTok}[1]{\textcolor[rgb]{0.49,0.56,0.16}{#1}}
\newcommand{\BaseNTok}[1]{\textcolor[rgb]{0.25,0.63,0.44}{#1}}
\newcommand{\BuiltInTok}[1]{\textcolor[rgb]{0.00,0.50,0.00}{#1}}
\newcommand{\CharTok}[1]{\textcolor[rgb]{0.25,0.44,0.63}{#1}}
\newcommand{\CommentTok}[1]{\textcolor[rgb]{0.38,0.63,0.69}{\textit{#1}}}
\newcommand{\CommentVarTok}[1]{\textcolor[rgb]{0.38,0.63,0.69}{\textbf{\textit{#1}}}}
\newcommand{\ConstantTok}[1]{\textcolor[rgb]{0.53,0.00,0.00}{#1}}
\newcommand{\ControlFlowTok}[1]{\textcolor[rgb]{0.00,0.44,0.13}{\textbf{#1}}}
\newcommand{\DataTypeTok}[1]{\textcolor[rgb]{0.56,0.13,0.00}{#1}}
\newcommand{\DecValTok}[1]{\textcolor[rgb]{0.25,0.63,0.44}{#1}}
\newcommand{\DocumentationTok}[1]{\textcolor[rgb]{0.73,0.13,0.13}{\textit{#1}}}
\newcommand{\ErrorTok}[1]{\textcolor[rgb]{1.00,0.00,0.00}{\textbf{#1}}}
\newcommand{\ExtensionTok}[1]{#1}
\newcommand{\FloatTok}[1]{\textcolor[rgb]{0.25,0.63,0.44}{#1}}
\newcommand{\FunctionTok}[1]{\textcolor[rgb]{0.02,0.16,0.49}{#1}}
\newcommand{\ImportTok}[1]{\textcolor[rgb]{0.00,0.50,0.00}{\textbf{#1}}}
\newcommand{\InformationTok}[1]{\textcolor[rgb]{0.38,0.63,0.69}{\textbf{\textit{#1}}}}
\newcommand{\KeywordTok}[1]{\textcolor[rgb]{0.00,0.44,0.13}{\textbf{#1}}}
\newcommand{\NormalTok}[1]{#1}
\newcommand{\OperatorTok}[1]{\textcolor[rgb]{0.40,0.40,0.40}{#1}}
\newcommand{\OtherTok}[1]{\textcolor[rgb]{0.00,0.44,0.13}{#1}}
\newcommand{\PreprocessorTok}[1]{\textcolor[rgb]{0.74,0.48,0.00}{#1}}
\newcommand{\RegionMarkerTok}[1]{#1}
\newcommand{\SpecialCharTok}[1]{\textcolor[rgb]{0.25,0.44,0.63}{#1}}
\newcommand{\SpecialStringTok}[1]{\textcolor[rgb]{0.73,0.40,0.53}{#1}}
\newcommand{\StringTok}[1]{\textcolor[rgb]{0.25,0.44,0.63}{#1}}
\newcommand{\VariableTok}[1]{\textcolor[rgb]{0.10,0.09,0.49}{#1}}
\newcommand{\VerbatimStringTok}[1]{\textcolor[rgb]{0.25,0.44,0.63}{#1}}
\newcommand{\WarningTok}[1]{\textcolor[rgb]{0.38,0.63,0.69}{\textbf{\textit{#1}}}}
\usepackage{longtable,booktabs,array}
\usepackage{calc} % for calculating minipage widths
% Correct order of tables after \paragraph or \subparagraph
\usepackage{etoolbox}
\makeatletter
\patchcmd\longtable{\par}{\if@noskipsec\mbox{}\fi\par}{}{}
\makeatother
% Allow footnotes in longtable head/foot
\IfFileExists{footnotehyper.sty}{\usepackage{footnotehyper}}{\usepackage{footnote}}
\makesavenoteenv{longtable}
\setlength{\emergencystretch}{3em} % prevent overfull lines
\providecommand{\tightlist}{%
  \setlength{\itemsep}{0pt}\setlength{\parskip}{0pt}}
\setcounter{secnumdepth}{-\maxdimen} % remove section numbering
\usepackage{amsmath,amssymb}
\usepackage{fvextra}
\fvset{breaklines=true,breakanywhere=true,fontsize=\scriptsize}
\ifLuaTeX
  \usepackage{selnolig}  % disable illegal ligatures
\fi
\usepackage{bookmark}
\IfFileExists{xurl.sty}{\usepackage{xurl}}{} % add URL line breaks if available
\urlstyle{same}
\hypersetup{
  pdftitle={Envelope contact on products of finitely many balanced graphs},
  pdfauthor={Xean research manuscript},
  hidelinks,
  pdfcreator={LaTeX via pandoc}}

\title{Envelope contact on products of finitely many balanced graphs}
\author{Xean research manuscript}
\date{4 October 2026}

\begin{document}
\maketitle

\subsection{Statement, coefficient convention, and
scope}\label{statement-coefficient-convention-and-scope}

Let \(d\geq2\), and let \(G_j\subset N_{j,\mathbb{R}}\) be nonempty
finite embedded rational pure graphs with strictly positive actual
\textbf{real} balanced edge weights. Use the direct product ambient
lattice, the actual saturated tangent lattices, and product weights on
\(X=\prod_{j=1}^dG_j\). Bounded edges, cycles, full lines, ends and
disconnected components are allowed. At a vertex use its actual
primitive outgoing vectors \(q_e\), and set \[
 K_v=\{c\geq0:\sum_e c_eq_e=0\}.
\] Rank means \(\dim\operatorname{span}_{\mathbb{R}}K_v\). Branched
means that some carrier junction has valence at least three; rational
bivalent subdivision does not count as branching. A graph presented
without vertices is treated by inserting rational bivalent points.

Real PA throughout means continuous real-affine pieces on a
\textbf{finite rational spatial complex}. This is the real scalar
extension of the literal PA convention on each fixed complex. Cycles,
including local testing cycles, have real weights. We retain the
source's genuine PCCs: rational endpoints, real nonnegative coefficients
summing to one, and their actual, possibly irrational, barycentre; on
some allowed rational structure a common face containing that centre is
a face of every endpoint coface. Raw PC means every such Jensen
inequality. Raw competitors are finite functions; their continuity is
never presumed.

Fix any prescribed continuous rational-domain real-PA actual-PC
\(\gamma\) positive at infinity: on every affine reference cell its
slope is positive on every nonzero recession vector. For bounded
continuous rational-domain real-PA \(u\), set \(f=\gamma+u\), and use
exactly the raw envelope \[
 P_\gamma(u)(x)=\sup\{h(x):h\in\mathrm{PC}(X),\ h\leq f,
                                      h\leq\gamma+C_h\}.
 \tag{E}
\] The growth restriction here is only an upper bound. PAPC tests the
divisor against every nonnegative real balanced local cycle, in each
dimension \(1,\ldots,d\), on every ordinary open subset. The operator is
the complete recursive zero-cycle \(P^d\cdot[X]\), with no division by
\(d!\).

\textbf{Finite-product classification theorem.} For this
fixed \((X,\gamma)\), every bounded continuous rational-domain real-PA
obstacle has a finite continuous rational-domain actual-PAPC raw
envelope in \(\gamma+O(1)\), a uniform decreasing PAPC regularization,
and \[
       \int_X(P_\gamma(u)-\gamma-u)\,
                  \mathrm{MA}_{\mathrm{poly},\mathbb{R}}(P_\gamma(u))=0
 \tag{C}
\] if and only if every local cone in every factor has rank one and
\textbf{at most one factor is branched}.

When a cone in one factor has higher rank, for each prescribed reference
we construct a compactly supported nonnegative obstacle with exact PAPC
raw envelope \(P\), with \(P=\gamma\) outside a compact subset of the
entire original product, and complete integral \[
 -d!\,(\delta\varepsilon M)\prod_{j=2}^d(2\delta b_jw_j)<0.
 \tag{H}
\] When all cones have rank one and two factors branch, the complete
integral is \[
 -\frac{d!}{2!}\,8\kappa^2ab\prod_{j=3}^d(2\kappa b_jw_j)<0.
 \tag{T}
\] Empty products in these formulas equal one. Every scalar is defined
below from actual weights and actual primitive directions, rather than
imposed as a mass assumption. Both failures have explicit regularization
\(P+1/n\).

This theorem classifies the specified finite-product family in the
declared real coefficient category. The original P05.2 classification
on arbitrary balanced carriers remains open. Literal integral-slope
coefficients, analytic realization, general singular-fan contact and
P06.2 are outside the theorem's scope. The complete original task,
frozen argument, supporting proofs, exact primary quotations and
independent review accompany this manuscript.

\subsection{Recursive intersection in every dimension, with saturated
quotients}\label{recursive-intersection-in-every-dimension-with-saturated-quotients}

Here and below
\(N_\sigma=N\cap\operatorname{span}_{\mathbb{R}}(\sigma-\sigma)\). For a
facet \(\tau\prec\sigma\), choose an integral lift of the primitive
inward generator of \(N_\sigma/N_\tau\). For a real PA \(h\) and a
balanced real \(k\)-cycle \(z\) define \[
 (h\cdot z)(\tau)=\sum_{\sigma\succ\tau}z(\sigma)h_\sigma^{\rm lin}(n_{\sigma/\tau})
       -h_\tau^{\rm lin}\left(\sum_{\sigma\succ\tau}z(\sigma)n_{\sigma/\tau}\right).
 \tag{D}
\] The sum in parentheses is tangential. Lift changes cancel by
agreement of linear parts on \(\tau\). Thus (D) is bilinear, local on
ordinary opens, and zero for ambient affine functions, including real
ones.

For refinement examine the minimal old face of a tested new
\((k-1)\)-face. If its dimension is above \(k\), all incident inherited
\(k\)-weights vanish. If it has dimension \(k\), exactly two pieces in
that inherited \(k\)-plane have possibly nonzero weights; they have
equal weight and linear part and opposite primitive quotient classes.
All other incident pieces have weight zero. Their contributions cancel.
If it has dimension \(k-1\), its tangent span is unchanged, and locally
there is one inherited incident \(k\)-piece for each old incident
\(k\)-face. The saturated lattices and normal classes are unchanged, so
the coefficient is the old coefficient. These cases also show inherited
balancing. They exhaust the possibilities since the old face contains
the new tested face. The proof applies to signed and real weights and at
every recursive stage. A newly inserted zero-face therefore inherits
zero top coefficient unless it was already an old zero-face.

For balancing of a new divisor, fix a \((k-2)\)-face \(\theta\) for
\(k\geq2\). Translate its germ, subtract an ambient real linear
extension of the common tangential linear part, and quotient by its
tangent lattice. Saturation of \(N_\theta\) follows from its definition.
The image \(N_\sigma/N_\theta\) is saturated in \(N/N_\theta\): if
\(m[n]\) belongs to it then \(mn+t\in N_\sigma\) for
\(t\in N_\theta\subset N_\sigma\), hence \(mn\in N_\sigma\), and
saturation implies \(n\in N_\sigma\). Consequently all primitive normals
are precisely those of the actual quotient. The incident cycle becomes a
balanced two-dimensional fan, possibly in a larger ambient quotient.
Rational homogeneous cuts make its two-cones pointed. Such cuts do not
change the divisor by the refinement argument.

For an unordered two-cone bounded by primitive rays \(r,s\), put
\(\chi_{rs}=[N_{C_{rs}}:\mathbb{Z}r+\mathbb{Z}s]\) and
\(\omega_{rs}=z(C_{rs})\). Extending \(r\) to a basis of the saturated
plane lattice gives the primitive normal class represented over the
reals by \(s/\chi_{rs}\). The discrepancy from an integral lift is
tangential and cancels in (D). Balancing at a ray and the first divisor
say \[
 \sum_s\frac{\omega_{rs}}{\chi_{rs}}s=t_rr,
 \qquad D_h(r)=\sum_s\frac{\omega_{rs}}{\chi_{rs}}h(s)-t_rh(r).
 \tag{F}
\] Here the function is homogeneous after subtraction. In
\(\sum_rD_h(r)r\), interchange \(r,s\) in the double first sum; it
becomes \(\sum_r h(r)\sum_s(\omega_{rs}/\chi_{rs})s\), exactly the
subtracted sum. Thus the new divisor is balanced at \(\theta\). For
\(k=1\), the new zero-cycle has no balancing condition. This proves
balancing in all dimensions.

For a second function also subtract its ambient tangential affine
extension. The full coefficient of its divisor on this balanced first
divisor is \[
 \sum_{\{r,s\}}\frac{\omega_{rs}}{\chi_{rs}}
                     [h(s)g(r)+h(r)g(s)]
                     -\sum_r t_rh(r)g(r).
 \tag{S}
\] Both terms are symmetric. In particular the tangential term is
retained. This proves mixed commutativity at every \((k-2)\)-face, hence
globally by refinement and locality. Adjacent transpositions give every
order of a mixed product. Since every intermediate cycle is balanced, an
ambient real affine term in any slot changes no product.

These arguments prove the real coefficient extension directly.
Alternatively, on a fixed rational complex both affine continuity and
cycle balancing are rational homogeneous linear systems: rational bases,
with denominators cleared, span their real solution spaces by integral
data. Bilinearity extends identities from those data. Positivity does
\textbf{not} follow from a signed-basis expansion; the full real
positivity tests are proved in the next section. The source
Allermann--Rau, arXiv:0709.3705v3, Definition 3.1 through Proposition
3.7, uses integral functions and integral cycle weights. Its full
selected body is included below for comparison. Its printed short
commutativity formula is not used in place of (S); no extension or
dropped tangential term is presumed.

Products of balanced cycles are balanced: their facets lower exactly one
factor, and their actual quotient normal belongs to that factor,
multiplied by the other factor weight. On a product cell,
\(N_{\sigma\times\eta}=N_\sigma\oplus N_\eta\), because the ambient
lattice is a direct sum. For a pulled-back function \(h\) from the first
factor, \[
             h\cdot(z\times z')=(h\cdot z)\times z'.
 \tag{P}
\] Indeed at \(\tau\times\eta\) both sums in (D) are \(z'(\eta)\) times
the first-factor sums, with primitive class \((n_{\sigma/\tau},0)\). At
\(\sigma\times\tau'\), normals \((0,n_{\eta/\tau'})\) give zero linear
parts and zero correction for \(h\). These are all facets. This remains
true for zero-dimensional factors: then every facet lowers another
factor, while the restricted function is constant, so it contributes
zero. This proves the zero-repeat-factor assertion, without a
transversality hypothesis or an index replacement. The same proof works
for every intermediate tensor cycle and for a two-dimensional first
factor.

For any continuous rational-domain real-PA PAPC \(P\), its successive
divisors are balanced and nonnegative; its top measure is the finite
zero-cycle on its domain vertices. Constants have zero linear parts.
Thus \(P+1/n\) decreases uniformly in \(\gamma+O(1)\) and has exactly
the same full measures. This constructs the regularized PA operator used
here directly; no general continuity theorem is required. There is no
diffuse mass or point at infinity.

For completeness, bounded reference shifts preserve the total top mass
in every dimension. Mixed commutativity gives \[
 P^d\cdot[X]-\gamma^d\cdot[X]
 =(P-\gamma)\cdot\sum_{i=0}^{d-1}P^i\gamma^{d-1-i}\cdot[X].
 \tag{M}
\] The inner cycle is a finite balanced real one-cycle. On a common
affine structure each bounded edge contributes opposite outgoing slopes
at its two endpoints, and each end has zero slope for the bounded affine
function \(P-\gamma\). Full lines can be split rationally and cancel.
The total difference is zero. This mass identity is separate from
contact.

\subsection{Finite PC cones, real-centre equivalence, and all local
positivity
tests}\label{finite-pc-cones-real-centre-equivalence-and-all-local-positivity-tests}

The elementary finite-polyhedron proof, all-centre/all-structure
finite-PC cone proof, exact active-face rational-box reference
decomposition, intrinsic graph criterion, and rank-one normalization
proofs are supplied in full immediately after the new theorem proof
below. They use neither the two-graph theorem nor the source's
continuity/locality conclusions. Their formulas (FC1)--(FC6), (G1) and
the vertex/recession formula (8) apply in \textbf{any} finite embedded
rational complex: none uses dimension two. We explicitly prove the
remaining dimension and coefficient points here.

If a continuous finite rational-domain real-PA \(h\) is actual PC, take
any nonnegative local real \(k\)-cycle on an ordinary open \(U\), for
any \(1\leq k\leq d\). Refine its global rational structure and that of
\(h\). At a tested facet choose rational
\(x\in U\cap\operatorname{relint}\tau\). Possible nonnegative incident
weight vectors balanced modulo the saturated tangent space form a
rational polyhedral cone. The supplied finite-generation proof gives
rational generators, so it suffices to test a rational vector
\(c_\sigma\). Write \(S=\sum c_\sigma>0\) and
\(T=\sum c_\sigma n_{\sigma/\tau}\), which is rational and tangential.
For small rational \(\epsilon>0\), the endpoints \[
              x_\sigma=x+\epsilon(n_{\sigma/\tau}-T/S)
\] are in their incident cells and \(U\). They are rational and have
barycentre \(x\) with coefficients \(c_\sigma/S\). PC Jensen multiplied
by \(S/\epsilon\) is exactly the nonnegativity of (D). The zero vector
gives zero. Linearity in weights proves every real test. This uses every
dimension and every ordinary open. The multiplier is \(S/\epsilon\).

Conversely, PAPC for continuous PA implies the finite tangent tests
(FC3). For a rational generator \((d_\sigma)\) of the sum-zero
tangent-family cone at a face \(\tau\), choose rational
\(y\in\operatorname{relint}\tau\). Discard zero directions and write
each nonzero rational vector \(d_\sigma=a_\sigma n_\sigma\), with
\(a_\sigma>0\) rational and \(n_\sigma\) a primitive integral vector on
its ray. Small initial segments are inside the corresponding cofaces by
(FC1). On a sufficiently small ordinary ball around \(y\), put weight
\(a_\sigma\) on those rays, adding weights when rays coincide. Extend
the segments to rational endpoints outside that ball. Their local
one-cycle is balanced because \(\sum d_\sigma=0\). Every other
subdivision point is bivalent, with equal opposite weights. This is an
allowed local cycle: a finite common ambient rational hyperplane
refinement includes each rational segment and \(y\), and all outer
endpoints are outside the chosen open. PAPC at \(y\) is
\(\sum h_\sigma^{\rm lin}(d_\sigma)\geq0\). If there are no nonzero
directions the assertion is zero. The finite-cone sufficiency proof
gives genuine PC at every real centre and every permitted structure. No
irrational centre is inserted as a rational vertex. This repairs
precisely that issue in a literal reading of the source's Lemma 3.4
proof.

For any number of graph factors, product-face tangent cones are
products: each projected direction has a sufficiently small feasible
positive step, and the minimum of finitely many such step sizes works
simultaneously. Given a sum-zero product tangent family, group its
projection by each factor coface. These grouped vectors belong to convex
factor tangent cones and sum to zero. Each factor PC gives a nonnegative
derivative sum. Adding gives (FC3) on the product. Its all-structure
sufficiency proves that a sum of pulled-back factor-PC functions is
actual PC. Zero projected displacements give zero; no assertion that
projected remote cofaces form a factor structure is needed. The argument
works for any finite number of factors, for all real centres, all
allowed cofaces and all real coefficients.

\subsection{Positive direction: component reduction with actual real
weights}\label{positive-direction-component-reduction-with-actual-real-weights}

At a rank-one vertex the strictly positive actual weight vector spans
the full kernel of its rational direction matrix: small positive and
negative perturbations of it realize every kernel vector in the span of
\(K_v\). Its sole relation has a positive primitive integer generator.
Incident actual weight ratios are consequently rational. On each
connected graph component, fix one actual edge weight as scale and
propagate ratios along finite paths. Given actual balancing makes these
ratios consistent around every cycle. Thus \(w_e=c r_e\) with \(c>0\)
real and \(r_e>0\) rational. At a bivalent vertex directions are
opposite primitive vectors and weights equal. A rank-one branching
vertex has no opposite pair, since that would be another nonnegative
relation.

Each unbranched component is a whole rational affine line: bivalent
propagation stays on one line, a finite endpoint would be one-valent and
unbalanced, and a closed bounded cycle cannot embed in that line.
Full-line components already have this property. When at most one factor
branches, every product component is \(G\times E\), where \(E\) is a
complete rational affine space of dimension \(m=d-1\) with saturated
lattice equal to the direct sum of the actual primitive line lattices.
If none branches, choose any one line as \(G\). Its actual cycle is a
single positive scalar \(c\prod_{j\ne G}w_j\) times the rationally
normalized cylinder cycle. The divisor formula is linear in its initial
cycle. Induction shows that this scalar multiplies the \textbf{entire}
recursively defined measure once; it is not raised to the \(d\)-th
power. PC competitors and envelopes do not depend on that scalar.

The complete multivariate cylinder reconstruction is supplied directly
below. It proves the following for connected rank-one rationally
weighted \(G\): for every nonseparable real rational-domain reference
and obstacle, the raw envelope on \(G\times E\) is finite
rational-domain PAPC in \(\gamma+O(1)\) and all noncontact top
coefficients vanish. We use its proof, not its previous review. The
proof uses the compact \textbf{multivariate} dual polyhedron \[
 D=\{p\in E^*:p(r_E)\leq\gamma^{\rm lin}(r)
                 \text{ for every vertical reference recession generator}\},
\] all projected spatial vertices, the greatest assignment in the full
graph incidence system (including every cycle), and rational
supergradient polytopes. Dual chamber vertices may be real. The spatial
maximization domains remain rational because the gradients of the node
assignments are rational, not because of rationality of those dual
vertices. These domains have zero-faces only over graph nodes.

At each possible noncontact atom \((v,z_0)\), its exposed dual face
\(C_x\) has rational affine-hull direction because \(z_0\) is rational
and all defining normals are rational. If \(\dim C_x<m\), the full germ
has a common rational lineality after affine subtraction, and the entire
top divisor is zero. If \(\dim C_x=m\), interior covectors uniquely
expose \(z_0\). Noncontact makes the corresponding node cap strict.
Raising only that LP coordinate shows its graph Laplacian is zero;
positivity and concavity then make every outgoing slope affine in \(p\).
The sole star relation factors those slopes through a real form \(A\)
and a rational map \(B\), giving the complete germ
\(A(s)+\operatorname{supp}_{C_x}(z+B(s))\). It factors after affine
subtraction through an \(m\)-dimensional complete rational fan. The full
recursive quotient vanishing proof below makes its \((m+1)\)-st divisor
zero. This covers all junctions, inserted bivalent nodes, tail bases,
fibre walls and open-edge strata. The supplied all-dimensional algebra
(D)--(S) validates each affine subtraction and every intermediate cycle.
There is no interpolation or separability substitute for this
multivariate argument.

Restoring each actual scalar proves contact on each component. There are
finitely many components; they are relatively open and closed, and a PCC
or local incident test stays in the component of its common face. The
finite structures, growth constants and measures combine componentwise.
This proves the positive implication of the theorem.

\subsection{A localization lemma for every fixed real
reference}\label{a-localization-lemma-for-every-fixed-real-reference}

Use the full finite-PC-cone and active-face box proof below to write
exactly \[
          \gamma=\sum_{a=1}^J\lambda_a\gamma_a,
                \lambda_a>0,\quad\sum_a\lambda_a=1,
 \tag{R}
\] where every \(\gamma_a\) has rational affine coefficients, is actual
PC on the same finite rational complex, and retains every strict
recession inequality. Choose rational product centre \(o\), ambient
lattice bases, and \(\rho(x)=\sum_j\|x_j-o_j\|_1\), with homogeneous
recession norm \(\rho_0\). For each \(a\), pointed-cell vertex/recession
decomposition gives rational \(m_a>0,Q_a\geq0\) such that
\(\gamma_a\geq m_a\rho-Q_a\) everywhere. Indeed choose \(m_a\) below all
positive rational ratios of reference recession slope to \(\rho_0\);
affinity bounds the vertex term below and the triangle inequality bounds
the norm above. Finitely many rational vertex bounds provide \(Q_a\). No
component or end is omitted. This also handles bounded cells.

Suppose a global rational-coefficient rational-domain actual-PC \(F\)
satisfies \(F\leq L\rho+A\) for rational \(L>0,A\). Choose rational
\(0<\delta_a<m_a/L\) and rational \(C_a\) sufficiently large that
\(C_a+\delta_a F(o)>\gamma_a(o)+2\). Define \[
 P_a=\max\{\gamma_a,C_a+\delta_a F\},\qquad
 P=\sum_a\lambda_a P_a,\qquad \delta=\sum_a\lambda_a\delta_a>0.
 \tag{L}
\] Maxima preserve all Jensen inequalities by selecting an active
function at the centre; positive sums also preserve them. Each maximum
has rational comparison walls because both operands have rational
coefficients. Their common rational refinement supports the real sum,
without any maximum with the original real \(\gamma\). The estimate \[
 \gamma_a-(C_a+\delta_aF)
 \geq(m_a-\delta_aL)\rho-Q_a-C_a-\delta_aA
\] proves \(P=\gamma\) outside a compact subset of the \textbf{entire}
product and \(P\geq\gamma\). Thus \(P\in\gamma+O(1)\) is actual PAPC
with the constant-shift regularization and retains each strict recession
slope. Its global compact localization has not modified the given
reference.

We can take a common function structure \(\Pi\) by finitely many ambient
rational hyperplanes extending all rational cell walls. Insert the
rational centre as an original product vertex (using bivalent
subdivisions in extra factors). Where all localizers dominate, the germ
is a finite homogeneous rational-domain germ plus a constant. Choose a
sufficiently small rational radius \(r_0>0\) such that \(\rho<r_0\) lies
in that germ, encounters no other factor endpoint, cap bend or
nonincident cell, and avoids every cutting hyperplane not through \(o\).
Every retained constraint is homogeneous about \(o\). A nonzero point
has a short radial contraction and expansion within its minimal face, so
is not a vertex. Hence \[
                    \Pi^{(0)}\cap\{\rho<r_0\}=\{o\}.
 \tag{V}
\] All inherited homogeneous cuts from the prescribed nonseparable
reference are retained. The norm ball is intersected with the carrier,
not identified with a product star.

\subsection{Higher-relation obstruction in any number of
factors}\label{higher-relation-obstruction-in-any-number-of-factors}

Suppose \(G_1\) has rank at least two at \(v\). Choose rational interior
points \(y_j\) of any edges of the other nonempty pure graphs, with
actual primitive tangents \(r_j\) and weights \(w_j>0\), for
\(2\leq j\leq d\). Their small germs are straight intervals even if the
graphs branch elsewhere. Set \(o=(v,y_2,\ldots,y_d)\). Let \(q_i,w_i\)
be the first star's actual directions and weights.

The rational pointed cone \(K_v\) has a rational extreme generator
\(c\ne0\) with nonempty proper support \(I\). If it were strictly
positive, an independent kernel vector would give two small opposite
perturbations of \(c\) in the cone, contradicting extremality. The full
kernel equals \(\operatorname{span}K_v\) by perturbing the strictly
positive actual weight vector. Put \(s_i=0\) on \(I\) and \(s_i=1\)
outside; then \[
       c\cdot s=0,\qquad b\cdot s\geq0\ (b\in K_v),
                  M=\sum_iw_is_i>0.
\] This is intrinsic graph data; it need not extend to an ambient convex
slope.

Set \(\rho_1=\|x_1-v\|_1\), \(\rho_j=\|x_j-y_j\|_1\). Along the first
ray \(\rho_1(v+tq_i)=a_it\), where \(a_i=\|q_i\|_1>0\). Choose rational
\(0<\varepsilon<\min_i a_i\), short rational star lengths \(L_i\) before
every existing endpoint or norm breakpoint, and rational
\(0<A<\min_iL_i(a_i-\varepsilon s_i)\). Replace the norm on this star by
\[
       F_1(v+tq_i)=\max\{a_it,A+\varepsilon s_it\},
\] and keep the norm elsewhere on every graph component. The comparison
points \(A/(a_i-\varepsilon s_i)\) are rational and before \(L_i\). At
the centre the slopes are \(\varepsilon s_i\); at each new edge bend
they increase to \(a_i\). Values agree at the centre and at the outer
gluing boundary. The intrinsic criterion (G1) proves global actual PC,
including opposite-pair relations and all real centres. It has rational
coefficients and domains, and \(0\leq F_1-\rho_1\leq A\).

Set \(F=F_1+\sum_{j=2}^d\rho_j\). The product tangent-family proof gives
global actual PC and PAPC in every dimension; \(\rho\leq F\leq\rho+A\).
Apply (L) with \(L=1\). On the central germ it gives \[
 P(v+tq_i,y_2+z_2r_2,\ldots,y_d+z_dr_d)
       =C+\delta\varepsilon s_it+\sum_{j=2}^d\delta b_j|z_j|,
       \qquad b_j=\|r_j\|_1>0,
 \tag{G}
\] where \(C=\sum_a\lambda_a(C_a+\delta_a A)\). The reference is
unchanged at infinity on all components. Choose the homogeneous rational
structure and radius of (V).

Choose rational \(B,T,D_j>0\), with \(B<T\), so small that
\(\max_i a_iT+\sum_{j=2}^d b_jD_j<r_0\). On this short product star put
\[
 \beta(v+tq_i,y_2+z_2r_2,\ldots)=
            (1-t/B-\sum_{j=2}^d|z_j|/D_j)_+,
 \tag{B}
\] and extend by zero. It vanishes before each outer boundary and has
the same value on the common first-vertex fibre. Thus it is continuous
compactly supported rational-domain rational-PA, equals one at \(o\),
and has support inside \(\rho<r_0\). Put \(u=P-\gamma+\beta\).

For any raw competitor \(q\leq P+\beta\), first fix a rational tuple
\(z\) with \(|z_j|\leq D_j\). The rational endpoints
\((v+Tq_i,y_2+z_2r_2,\ldots)\), \(i\in I\), with coefficients
\(c_i/\sum c_i\), form a genuine PCC centred on the entire apex fibre:
the supported relation cancels the first displacement and all other
displacements are zero. A permitted coarse rational product structure
has the closed fibre face over \(v\) in each incident strip; refine
extra sign coordinates if desired. The ceiling at these endpoints equals
\(P\) because \(T>B\). Equation (G) and \(c\cdot s=0\) give
\(q(v,z)\leq P(v,z)\). For any rational \(t\in[0,B]\), combine the
rational pinned apex and the endpoint at \(T\) in the same incident
strip; \(P\) is affine in \(t\). Hence \(q\leq P\) there. Outside the
bump, the ceiling gives the same bound directly. We have bounded every
rational point of the whole product, not just the central point.

For a real point with positive bump, choose its minimal face of the
local rational ray/sign-orthant structure, with \(t=0\) and all zero
\(z_j\) as faces. Arbitrarily small rational boxes in rational affine
coordinates enclose it within that face and the germ. Their real
barycentric coefficients give a same-face PCC. Every rational endpoint
satisfies the preceding bound, and (G) is affine on that face, so
\(q(x)\leq P(x)\) exactly. Zero-faces are rational; all boundary cases
with zero displacements are thus included. Elsewhere \(\beta=0\).
Conversely \(P\) is an actual-PC upper-growth competitor below
\(P+\beta\). Therefore (E) equals \(P\) on all of \(X\), without
assuming continuity or lower growth for \(q\).

The local cycle is \(Z_1\times\cdots\times Z_d\), where \(Z_1\) is the
actual star and \(Z_j\), \(j\geq2\), is the whole primitive line germ of
weight \(w_j\). Write \(h_1(tq_i)=\delta\varepsilon s_it\) and
\(h_j(z_jr_j)=\delta b_j|z_j|\). Their one-factor divisors are exactly
\[
 h_1\cdot Z_1=(\delta\varepsilon M)[0],\qquad
 h_j\cdot Z_j=(2\delta b_jw_j)[0].
 \tag{A1}
\] In the first formula the primitive outgoing directions are \(q_i\),
and in the second they are \(r_j,-r_j\). The tangential correction at
the point is zero; these are the complete one-cycles' divisors.

At every product stratum the saturated lattice is the direct sum of the
remaining primitive factor-line lattices. Lowering one factor therefore
uses its actual primitive vector, with index one in that split quotient.
This does not require the whole first-star span to be generated
saturately by its rays. For each \(0\leq r\leq d\), (P) and
multilinearity give the complete intermediate cycle \[
 (\sum_jh_j)^r\cdot\prod_jZ_j
 =r!\sum_{\substack{I'\subset\{1,\ldots,d\}\\|I'|=r}}
       \prod_{j\in I'}(h_j\cdot Z_j)\times\prod_{j\notin I'}Z_j.
 \tag{A2}
\] This follows by induction: words with a repeated factor give zero on
its zero-cycle, while every ordering of distinct factors has the same
tensor weights and primitive normals by (P); each subset has \(r!\) such
words. The formula lists every intermediate product stratum, not only
its eventual atom. All cycles are balanced and nonnegative. At \(r=d\)
it gives exactly the positive mass in (H). Additional homogeneous cuts
preserve it by (D)'s refinement proof at \textbf{all} stages. This
proves the coefficient, including every order and the factorial.

\subsection{Two rank-one branched factors: fibrewise pinning and all
stages}\label{two-rank-one-branched-factors-fibrewise-pinning-and-all-stages}

Suppose every cone has rank one but \(G_1,G_2\) branch at \(v_1,v_2\).
At their stars write the actual weights as \(a\ell_i\) and \(b m_j\),
where \(\ell,m\) are primitive positive integer relations and \(a,b>0\)
are actual real scales. Choose three nonempty groups \(I_k,J_k\), with
\(S_k=\sum_{i\in I_k}\ell_i\), \(T_k=\sum_{j\in J_k}m_j\), and tripod
vectors \(\nu_1=(1,0),\nu_2=(0,1),\nu_3=(-1,-1)\). The assignments
\(L_1(q_i)=\nu_k/S_k\), \(L_2(r_j)=\nu_k/T_k\) annihilate their
\textbf{entire} kernels. They therefore extend to rational ambient maps
by rational linear algebra, without altering the original star lattices.
Set \[
 H(x_1,x_2)=\|L_1(x_1-v_1)-L_2(x_2-v_2)\|_1.
\] It is ambient convex rational PA. Choose rational edge-interior
\(y_j\) in every remaining factor, with primitive \(r_j\), actual weight
\(w_j\), and \(b_j=\|r_j\|_1\). Let \(F=H+\sum_{j=3}^d\rho_j\). Ambient
map matrix bounds give \(F\leq L\rho\) for some rational \(L>0\). Apply
(L), denoting its positive sum scale by \(\kappa\). In a full product
germ around \(o=(v_1,v_2,y_3,\ldots,y_d)\), \[
                    P=C+\kappa H+\sum_{j=3}^d\kappa b_j|z_j|.
 \tag{J}
\] Choose a homogeneous structure and \(r_0\) satisfying (V). This
preserves \(\gamma\) outside a compact set on the entire original
product.

Here is the fibrewise repair of single-point pinning. Let
\(n(x_1,x_2)=\|x_1-v_1\|_1+\|x_2-v_2\|_1\). Choose rational \(\tau>0\)
so all finitely many endpoints \[
 e_{ij}=(v_1+\tau S_kq_i,v_2+\tau T_kr_j),\quad
 \alpha_{ij}=\ell_i m_j/(3S_kT_k),\quad i\in I_k,j\in J_k
 \tag{Q}
\] are in their original short cofaces and \(n(e_{ij})<r_0/4\). They are
nonzero, rational, and have \(H(e_{ij})=0\). Coefficients sum to one and
both projected barycentres are the star origins, by the original
relations. Choose rational extra-coordinate widths \(D_j>0\) with
\(\sum_{j=3}^db_jD_j<r_0/4\), and rational \(R\in(0,r_0/4)\). Choose
rational \(B>0\) smaller than \(R\) and every \(n(e_{ij})\). Put on the
whole product \[
        \beta=(1-n/B-\sum_{j=3}^d\rho_j/(b_jD_j))_+,
                    u=P-\gamma+\beta.
 \tag{K}
\] For \(d=2\) the sum is empty. Its support is contained in the chosen
local product germ: positive bump implies \(n<B\) and \(|z_j|<D_j\),
hence \(\rho<r_0/2\). This is a continuous compact rational norm tent,
with rational walls and coefficients and \(\beta(o)=1\).

Fix any rational tuple \(z\), \(|z_j|\leq D_j\). Tensor (Q) with those
fixed extra coordinates. It is a genuine PCC whose common face is the
whole bad-pair apex fibre, on the original product structure subdivided
rationally at the extra points. Its endpoints are in (J), and their bump
is zero since \(n(e_{ij})>B\). The ceiling bound and \(H(e_{ij})=0\) pin
\[
                  q(v_1,v_2,z)\leq C+\sum_{j=3}^d\kappa b_j|z_j|
                         =P(v_1,v_2,z).
 \tag{AP}
\] This pins the full apex fibre, not merely \(o\).

For a rational bad-pair point \(x\ne(v_1,v_2)\) with \(n(x)<B\), choose
one closed homogeneous rational \(H\)-affine cone containing it in the
product of the two short stars. Its closure contains the apex. Define
the rational outer point \(y=(R/n(x))x\) in displacement coordinates.
Then \(n(y)=R>B\), and \(y\), together with each fixed \(z\) under
consideration, is in the exact germ (J): \(R+\sum b_jD_j<r_0/2\). Both
points lie before all original bounded endpoints by the choice of germ.
The tent vanishes there. The point \((x,z)\) is the same-coface PCC of
\((v_1,v_2,z)\) and \((y,z)\), with coefficients \(1-n(x)/R,n(x)/R\). A
rational local subdivision retaining the closed \(H\)-affine cone and
the extra sign faces, extended by ambient rational cuts, is permitted.
We may take that cone cell as common face for a same-cell PCC; both
endpoints lie in it. Affinity of (J), (AP), and the outer ceiling give
\(q(x,z)\leq P(x,z)\). If the bad-pair displacement is zero use (AP),
and if the bump is zero use the ceiling directly. This covers every
rational point, including all cone boundaries and zero extra
displacements.

At a real point with positive bump use a small rational enclosing box or
simplex in its minimal face of the homogeneous \(H\)-affine/sign
structure. It stays in the germ; all endpoints are rational and already
pinned, and \(P\) is affine on that face. The real-centre same-face PCC
gives \(q\leq P\) exactly. At zero-faces the point is rational.
Elsewhere the bump vanishes. Since \(P\) itself is an allowed global
competitor, this proves the true equality \(P_\gamma(u)=P\) on the
entire product. It requires neither raw continuity nor a stronger growth
condition nor irrational endpoints.

We supply the complete original bad-pair mass calculation below, without
a tripod pushforward. In quadrant \((sq_i,tr_j)\), write
\(q'_i=1/S_k,p'_j=1/T_l,c_k=(1,1,2)_k\). Then \(H=c_kq'_is+c_lp'_jt\)
off-group and \(H=c_k|q'_is-p'_jt|\) in-group. Actual saturated planes
have primitive axes \(a_i=(q_i,0),b_j=(0,r_j)\). For a common
denominator \(D\), \(Q=Dq'_i,U=Dp'_j,g=\gcd(Q,U)\), the kink ray is the
primitive \(d_{ij}=(U/g)a_i+(Q/g)b_j\), with cone indices \(Q/g,U/g\).
Its first weight is \(2c_kab\ell_i m_jg/D\) and its value is zero. At an
axis the in-group quotient representative is \(b_j+(p'_j/q'_i)a_i\). The
full tangential correction gives first axis weights
\(ab\ell_i(4-2c_k)\), \(ab m_j(4-2c_l)\). Axis plus kink first
contributions have projections \(4ab\ell_iq_i\), \(4ab m_jr_j\), so
their first cycle is balanced. The second coefficient is \(4ab\) from
each axis family and zero from all kinks, namely \(8ab\). All normals
and indices are directly derived in the included detailed calculation.
No shortened version of (S) is used.

Let \(Z\) denote this actual bad-pair two-cycle, \(Z_j\) the extra
primitive line germs, \(h_0=\kappa H\), and \(h_j=\kappa b_j|z_j|\). For
\textbf{every} \(0\leq r\leq d\), the entire intermediate cycle is \[
 (h_0+\sum_{j=3}^dh_j)^r\cdot(Z\times\prod_{j=3}^dZ_j)
 =\sum_{\substack{0\leq n_0\leq2,\ n_j\in\{0,1\}\\n_0+\sum n_j=r}}
       \frac{r!}{n_0!}\,(h_0^{n_0}\cdot Z)
                              \times\prod_{j=3}^d(h_j^{n_j}\cdot Z_j).
 \tag{A3}
\] The factorial denominator accounts for the two identical bad-pair
slots. Prove this by expansion into words and (P): lowering an extra
factor twice is zero, lowering the bad pair a third time is zero on its
zero-cycle, and every surviving order is the displayed tensor with
inherited weights. Product saturated lattices split; the bad-pair
primitive kink quotient normals and indices remain exactly \(Q/g,U/g\)
when extra line lattices are added or removed. Other quotient normals
are the primitive extra tangents. Thus (A3) retains all axis, kink,
multiple-order and zero-repeat contributions. At \(r=d\) only
\(n_0=2,n_j=1\) survives, giving exactly the mass in (T). This derives
its coefficient from all stages rather than guessing a tensor mass.

\subsection{Complete integrals and exhaustion of the
criterion}\label{complete-integrals-and-exhaustion-of-the-criterion}

In both failures the common old structure \(\Pi\) has only \(o\) as a
zero-face where the bump can be positive, by (V). The actual envelope
has exact gap \(-\beta\), equal to \(-1\) there and zero at every other
old measure vertex. Its measure is the finite nonnegative complete
\(d\)-stage PA zero-cycle. Refining for the obstacle can introduce
vertices with positive bump, but at every stage refinement gives the
inherited cycle; any zero-face that was not an old zero-face inherits
zero top coefficient. There is no mass on open positive-dimensional
strata, no diffuse term, and no point at infinity. The full measures of
\(P+1/n\) are identical, so regularization adds no escaped or boundary
mass. Consequently the \textbf{entire} contact integrals are exactly (H)
and (T), not merely negative contributions at an isolated atom. All
remote components, cycles, bounded edges and ends have zero gap at their
possible measure atoms. Formula (M) also preserves total reference mass.

If a cone has higher rank the first construction works with arbitrary
other factor graphs, by their rational edge interiors. If every cone has
rank one but two factors branch, the second construction applies with
all remaining factors tested at edge interiors. Otherwise at most one
factor branches and the positive cylinder reduction applies. These
alternatives exhaust the proposed geometric criterion, including
failures occurring in different components of different factors. The
localizations are global, so any failed component tuple suffices against
the originally prescribed reference.

The exact finite-product theorem has passed an independent Xean
review. All supporting mathematical proofs are retained below. Source
quotations and comparison dossiers accompany the complete frozen
argument. Same-model review, finite arithmetic, compilation and file
integrity are not formal verification or priority certification.

\section{Complete elementary supporting proofs, reconstructed in this
argument}\label{complete-elementary-supporting-proofs-reconstructed-in-this-argument}

These proof bodies are included directly rather than invoked as prior
results. Formula references are internal to each reproduced proof unit.
The arbitrary-dimensional algebra and scalar normalization required by
their applications are proved in the main text. In the rank-one coface
discussion, replacing the complete one-dimensional affine factor by E
does not change the graph projection; all affine E terms cancel by the
full vector barycentre.

\subsection{Elementary finite-dimensional facts used in the
reconstruction}\label{elementary-finite-dimensional-facts-used-in-the-reconstruction}

A polyhedral cone has finitely many generators; for rational
inequalities these can be rational. Remove its lineality (a rational
subspace for rational data) and work in the pointed quotient. The sum of
its nonnegative defining forms is strictly positive away from zero:
otherwise all forms vanish on a nonzero vector and its negative. The
section of value one is bounded, since an unbounded sequence, normalized
to the unit sphere, would yield a nonzero cone vector of value zero. A
bounded polyhedron is the convex hull of its vertices: in the minimal
face of any nonvertex, a line ends on smaller faces, and induction
reduces dimension. Its vertices are the finitely many full-rank
active-system solutions, rational for rational data. Lift back and add
both signs of a lineality basis. Homogenizing a nonempty pointed
rational polyhedron therefore proves \[
 \sigma=\operatorname{conv}\{v_1,\ldots,v_l\}
                    +\operatorname{cone}\{r_1,\ldots,r_m\},     \tag{8}
\] with rational vertices and recession vectors and at least one vertex:
the height-one slice needs a positive-height extreme ray, whose
extremality gives an original vertex. The homogenized cone is pointed
because any line has height zero and would give original lineality. The
same proof holds over real data without rational coordinates.

Eliminating a variable from finitely many nonstrict linear inequalities
compares every lower bound with every upper bound and retains
variable-independent conditions. A one-sided family imposes no further
conditions. Repetition proves that projections of polyhedra are
polyhedra; variable coefficients stay rational if originally rational,
even with real constants. It also proves closedness of finite
convex-hull-plus-cone sets by projecting their coefficient
representation.

Needed separations can be proved by nearest points. A point outside a
nonempty closed convex set has a nearest point after restriction to a
sufficiently large compact ball. Differentiating squared distance along
segments gives a separating normal. A compact convex set disjoint from a
closed convex set has positive attained distance, with the same compact
restriction for the second set; differentiating in both sets strictly
separates them. If the latter is a cone, a finite supremum of the
separating functional on it is zero. This proves the cone/polytope
separations used here.

If the pointed cone quotient is zero-dimensional, only its zero vector
and the lifted lineality basis are needed. The elementary fact about a
bounded polyhedron used above also covers its zero-dimensional case.
Pointedness ensures the homogenized cone has no lineality; its
zero-height extreme rays give recession generators and its
positive-height extreme rays give exactly the normalized vertices. This
proves (8) without an unstated decomposition theorem.

A finite rationally convex function on \(\mathbb{Q}\) has a unique
continuous convex extension on any interval interior. Rational endpoint
values bound it above on a rational compact interval. Reflection about a
rational centre bounds it below on a smaller interval. Bounds
\(L\leq F\leq M\) on a larger rational interval imply a Lipschitz bound
on a smaller one: extend the rational segment from \(y\) through \(z\)
by a rational factor \(R/|z-y|\) staying in the larger interval and use
Jensen, then reverse \(y,z\). Density gives the completion and
convexity. A finite convex function on \(\mathbb{R}\) is continuous by
this argument, and has a touching affine support at each point: its left
and right secant limits are finite, ordered, and any slope between them
gives the support.

For \(b(p)=\min_j(c_j+g_jp)\), the supergradient interval at \(p_0\) is
the convex hull of the active \(g_j\). Its directional derivative is the
minimum of their products with the direction. Separating a candidate
outside that interval disproves its supergradient property. Expanding
sums of finite minima proves that supergradients of nonnegative sums are
the corresponding weighted Minkowski sums, also when coefficients
vanish. The outward normal cone of a polytope is generated by active
inequality normals: an alleged normal outside their finitely generated
closed cone can be separated to give a direction satisfying the active
inequalities and increasing the alleged normal; a sufficiently small
step also satisfies inactive inequalities, a contradiction. These proofs
retain rational normal generators whenever inequality coefficients are
rational, even if constants are real.

\subsection{A finite rational cone for actual PC on a fixed
complex}\label{a-finite-rational-cone-for-actual-pc-on-a-fixed-complex}

This lemma concerns a \textbf{fixed} finite embedded rational complex
\(\Sigma\). It does not assert regularity on its carrier. Let
\(A(\Sigma)\) be the real vector space of continuous functions affine on
each cell, with real affine coefficients. Choose rational affine
coordinates on each cell's affine span. Agreement on each common face is
a finite rational homogeneous linear system: test agreement on a
rational affine basis of that face. These tests suffice because the
difference is affine. Thus \(A(\Sigma)\) has a rational basis, including
when cells have lineality or are lower dimensional. Clearing slope
denominators in a rational basis element puts it in the literal PA
class. This proves that the declared function space is its real scalar
extension.

For every nonempty face \(\tau\), every coface \(\sigma\succeq\tau\),
and \(y\in\operatorname{relint}(\tau)\), define \[
 T_{\tau\sigma}=\operatorname{cone}(\sigma-y).                 \tag{FC1}
\] Here cone means nonnegative scalar multiples; this set is already a
closed convex polyhedral cone. To prove this, write \(\sigma\) using
rational affine equalities and inequalities \(a_k(z)+b_k\geq0\). At
\(y\), an inequality is active if and only if it vanishes identically on
\(\tau\): a nonnegative affine function on \(\tau\) cannot vanish at a
relative-interior point unless its restriction is zero, by taking short
steps in both tangent directions. A vector \(d\) is in (FC1) exactly
when it satisfies the homogeneous equalities and \(a_k(d)\geq0\) for all
active inequalities. Necessity follows by substitution. For sufficiency,
every inactive inequality has positive slack at \(y\); there are
finitely many, so a sufficiently small positive step \(y+t d\) satisfies
all of them, as well as the active constraints. Therefore
\(y+t d\in\sigma\). This also proves that each such direction stays in
\(\sigma\) for every sufficiently small positive step. The same active
constraints hold at every real point in \(\operatorname{relint}(\tau)\),
proving independence of the chosen centre, rationality, and inclusion of
the whole tangent space of \(\tau\) as lineality. No boundedness,
pointedness, or full dimension was used.

Form the cone of vector families, indexed by \textbf{all} cofaces
including \(\tau\), \[
 B_\tau=\{(d_\sigma):d_\sigma\in T_{\tau\sigma},\quad
                              \sum_{\sigma\succeq\tau}d_\sigma=0\}.
                                                               \tag{FC2}
\] It is a rational polyhedral cone and has finitely many rational
generators, with both signs of a rational lineality basis retained. The
finite-generation proof above applies to the product of the tangent
cones and the rational sum-zero equations. Impose, for every generator
and every \(\tau\), \[
                \sum_{\sigma\succeq\tau}h_\sigma^{\rm lin}(d_\sigma)
                                  \geq0.                       \tag{FC3}
\] These are rational homogeneous linear inequalities in \(A(\Sigma)\).

\textbf{Claim.} Their solution cone \(C(\Sigma)\) equals
\(A(\Sigma)\cap\mathrm{PC}(X)\), with actual PC tested against all
allowed rational structures, all real centres, and remote rational
endpoints.

For necessity choose a rational \(y\in\operatorname{relint}(\tau)\).
Such a point exists: rational affine coordinates identify the relative
interior with a nonempty open subset cut out by rational constraints,
and rational points are dense in its affine span. Given a rational
generator of \(B_\tau\), let \(m\) be the number of cofaces. A
sufficiently small rational \(t>0\) puts all rational endpoints
\(y+mt d_\sigma\) in their cofaces. Their equal coefficients \(1/m\)
give a genuine \(\Sigma\)-PCC centred at \(y\). Since every piece agrees
at \(y\), Jensen is precisely (FC3) after division by \(t\). The zero
family causes no difficulty. Linear combination extends the inequalities
from generators to every real family in \(B_\tau\).

For sufficiency first take any straight segment contained in the
carrier, with endpoints allowed to be real. Each cell intersects it in
an interval, so a finite subdivision makes \(h\) affine on each
subinterval. At an interior breakpoint
\(x\in\operatorname{relint}(\tau)\), let \(d\) be the forward direction.
Its incoming \(-d\) and outgoing \(d\) lie in the tangent cones of their
respective initial cells, which contain \(\tau\) as a face. Put them in
the corresponding slots of (FC2). If the slots coincide, sum them in
that slot; both derivatives then come from the same affine piece and
agree. In all cases (FC3) gives \(h'_+(x;d)-h'_-(x;d)\geq0\). A
continuous PA function with nondecreasing successive slopes is convex:
the secant slope over an interval is a weighted average of its piece
slopes, so left secants do not exceed right secants. Thus \(h\) is
convex on every contained straight segment.

Now take any genuine PCC \((x,(x_i),(\alpha_i))\) on \textbf{another}
permitted structure. Each segment \([x,x_i]\) lies in its permitted
convex coface and hence in \(X\). If \(d_i=x_i-x\ne0\), its initial
short segment lies in some \(\Sigma\)-cell containing \(x\); finiteness
of the segment subdivision ensures such a cell exists. The minimal
\(\Sigma\)-face \(\tau\) of \(x\) is a face of that cell, and
\(d_i\in T_{\tau\sigma_i}\) by (FC1). Continuity and affine agreement
make the initial derivative unambiguous. Convexity along the segment
gives \(h(x_i)-h(x)\geq h'_{x}(d_i)\). Group
\(D_\sigma=\sum_{i:\sigma_i=\sigma}\alpha_i d_i\). Each group stays in
its convex tangent cone, and \(\sum D_\sigma=\sum\alpha_i(x_i-x)=0\).
Therefore (FC3) gives \[
 \sum_i\alpha_i(h(x_i)-h(x))\geq
 \sum_\sigma h_\sigma^{\rm lin}(D_\sigma)\geq0.              \tag{FC4}
\] Zero displacements and zero coefficients contribute zero. This proves
the claim at \textbf{every real centre}, with no insertion of irrational
endpoints or cycles into the source definitions.

\subsection{Exact positive-growth decomposition of the prescribed real
reference}\label{exact-positive-growth-decomposition-of-the-prescribed-real-reference}

Choose a common finite rational affine structure \(\Sigma\) for the
product and the prescribed real \(\gamma\). Since \(\gamma\) is positive
at infinity, each cell is pointed: nonzero lineality would put both
\(r\) and \(-r\) in its recession cone, requiring opposite affine slopes
both positive. By (8), each recession cone has finitely many nonzero
rational generators. Its positive-growth conditions are exactly the
finite strict inequalities \[
       h_\sigma^{\rm lin}(r)>0\quad
       \text{for all chosen nonzero recession generators}.     \tag{FC5}
\] A nonzero recession vector is a nonnegative combination with at least
one positive coefficient, so these inequalities suffice as well as being
necessary. Bounded cells impose none.

Write \(C(\Sigma)\) by its finite inequalities \(L_i\geq0\) in the
rational space \(A(\Sigma)\), and let \(I=\{i:L_i(\gamma)=0\}\),
\(S=\bigcap_{i\in I}\ker L_i\). Then \(F=C(\Sigma)\cap S\) is the
minimal face containing \(\gamma\). Indeed it is a face by nonnegativity
of the \(L_i\). In \(S\), all other inequalities are strictly positive
at \(\gamma\), so a small relative ball about \(\gamma\) lies in \(F\).
Hence \(\operatorname{span}F=S\) and
\(\gamma\in\operatorname{relint}F\). Any face containing \(\gamma\)
contains that ball: for \(z\in F\) a short step from \(\gamma\) opposite
\(z-\gamma\) stays in the ball, expressing \(\gamma\) as a strict convex
combination of \(z\) and another point of \(F\); the defining face
property then includes \(z\). This proves minimality without a face
theorem. The subspace \(S\) is rational, even when the cone has
lineality.

Choose a rational basis of \(S\), and use it as coordinates. If its
dimension is \(d>0\), choose rational endpoints \(a_k<t_k<b_k\) about
the real coordinates \(t_k\) of \(\gamma\), so close that their entire
box lies in the preceding relative ball and satisfies every strict
inequality (FC5). Density of the rationals and finiteness of all
constraints give such a box. All its \(2^d\) vertices \(\gamma_j\) have
rational affine coefficients, belong to \(C(\Sigma)\), and retain every
strict recession slope. Put \(\theta_k=(t_k-a_k)/(b_k-a_k)\in(0,1)\).
The product weights obtained by choosing \(\theta_k\) at an upper
endpoint and \(1-\theta_k\) at a lower endpoint are strictly positive
and sum to one. Their coordinate averages are exactly \(t_k\). Thus,
\textbf{as functions on the entire unchanged carrier}, \[
                 \gamma=\sum_{j=1}^{J}\lambda_j\gamma_j,
                 \lambda_j>0,\qquad\sum_j\lambda_j=1.        \tag{FC6}
\] Each \(\gamma_j\) is actual PC against every structure by (FC4), not
just feasible in a nodal relaxation. If \(d=0\), then \(\gamma=0\) in
this space; use the single rational reference \(0\) with weight one. In
that case (FC5), if nonempty, would exclude \(0\), so this case can
occur only when there are no nonzero recession directions. This handles
the degenerate case without claiming growth where it is impossible. No
approximation, reference replacement, or lost strict recession
inequality occurs in (FC6).

\subsection{Intrinsic graph convexity and product projection, without
rank
assumptions}\label{intrinsic-graph-convexity-and-product-projection-without-rank-assumptions}

This section supplies the graph and projection lemmas used in the new
higher-relation case. It does not invoke the existing one-dimensional
classification.

For a continuous finite PA function \(h\) on a finite rational graph
structure, actual PC is equivalent to edgewise convexity and \[
             c\mathbin{\cdot}s_v(h)\geq0\quad(c\in K_v)       \tag{G1}
\] at every vertex, including function breakpoints. To prove this
equivalence, use (FC1)--(FC4) on that graph structure. At an edge
interior the only coface tangent cone is its full line, so a tangent
family of sum zero has zero affine derivative sum. At a vertex the
one-face tangent cones are the nonnegative primitive rays; every family
is \((t_eq_e)_e\), \(t_e\geq0\), with \(t\in K_v\), and the zero-face
contributes zero. Thus its tests are exactly (G1). At a bivalent affine
breakpoint \(K_v\) is generated by \((1,1)\), and (G1) says the right
derivative is at least the left derivative. These inequalities imply
edgewise convexity by monotone successive slopes. Conversely all such
tests imply (FC3), which proves PC for every real centre, remote
rational endpoint and allowed coface structure. At original vertices one
need not replace the cone \(K_v\) by the single actual weight relation.
All its real members are tested.

For completeness this also proves convexity on any contained straight
segment passing through a higher-rank branching vertex: its incoming and
outgoing opposite primitive directions give a member of \(K_v\), so its
successive slopes cannot decrease. No ambient affine extension of
\(s_v(h)\) is asserted.

If \(h_i\) are continuous actual-PC rational-domain PA functions on the
two graphs, then \[
                  h(x,y)=h_1(x)+h_2(y)                       \tag{G2}
\] is actual PC on their full product. Here is the precise projection
argument, which does not assume that arbitrary projected cofaces
themselves form a graph structure. Choose fixed rational function
structures \(\Sigma_i\). For a product face \(\tau_1\times\tau_2\) and a
product coface \(\sigma_1\times\sigma_2\), the tangent cone (FC1) is the
direct product \(T_{\tau_1\sigma_1}\times T_{\tau_2\sigma_2}\). Indeed
each projected direction admits a sufficiently small positive feasible
step, so both admit a common small step; the reverse inclusion follows
by projection.

For a product tangent family of sum zero, project its vectors and group
them over \(\sigma_1\). The grouped first-factor vectors stay in their
convex tangent cones and sum to zero. Factor PC and (FC3) give a
nonnegative first-factor derivative sum. Grouping over \(\sigma_2\)
gives a nonnegative second-factor sum. Their sum is the product test.
Consequently (FC4) proves (G2) against every genuine product PCC on
every other allowed structure. More explicitly, project the segments of
such a PCC, use factor convexity along each projected segment, and group
their initial directions in the fixed factor tangent cones. The
projected barycentre equations give precisely these sum-zero families. A
zero projected displacement has derivative zero; all projected distances
and all real coefficients are retained. This explains why projection is
valid even when a projected distant segment crosses a vertex. It does
not make an unsupported assertion about projected faces.

By the all-ordinary-open PC-to-PAPC proof above, (G2) and every PC
function constructed below are PAPC against every nonnegative real local
one- or two-cycle. Product balancing and the actual saturated
primitive-normal recursion, refinement and locality have already been
reconstructed.

Every balanced graph has a strictly positive vector \(w\in K_v\). Hence
\(\operatorname{span}K_v=\ker Q_v\), where \(Q_v\) has columns \(q_e\):
any kernel vector can be added to and subtracted from \(w\) after a
sufficiently small scaling while preserving nonnegativity. This proves
the equality and shows that a rank-one cone is precisely the sole
positive kernel relation. It also shows that higher rank cannot be
hidden by irrational actual weights. Rational bivalent subdivisions have
rank one and do not change any original relation cone or the carrier
branching property. These are local geometric conditions independent of
the chosen finite refinement.

\subsection{Rank-one supporting graph geometry, real weights, and
genuine
cofaces}\label{rank-one-supporting-graph-geometry-real-weights-and-genuine-cofaces}

From this point through the two-branch reconstruction, assume all local
relation cones in both graphs have rank one. This is a supporting case,
not a new hypothesis in the broader theorem. At a vertex the full kernel
of the rational direction matrix is one-dimensional. Otherwise perturb
its strictly positive actual weight vector by a sufficiently small
nonproportional kernel vector; positivity persists, contradicting the
rank-one condition. The kernel therefore has a primitive positive
integer generator \(\lambda_v\), and the actual incident vector is
\(a_v\lambda_v\), \(a_v>0\). At a bivalent vertex the directions are
opposite primitive vectors and weights equal. At a branched vertex there
is no opposite pair, since that pair would give another nonnegative
relation. No one-valent vertex can balance.

On a connected graph all edge weights are a single positive real scalar
times positive rational weights. Fix one edge weight. At each adjacent
vertex the ratios of incident weights equal ratios of positive integers
in \(\lambda_v\); propagate along paths through the finite graph.
Existing actual weights ensure consistency around every cycle. This
yields rational normalized weights \(r_e\) with \(w_e=c r_e\). Full-line
components have one arbitrary positive scalar weight. Rational
subdivision preserves these statements. This reduction is essential: it
makes the incidence elimination below have rational variable
coefficients, without supposing the given real scale rational.

Every connected unbranched component is a complete rational affine line.
Propagating through bivalent vertices keeps every edge on the same
affine line with equal weight. The component is a connected finite
closed union of intervals/rays on that line. Any finite endpoint of its
union would be one-valent, contrary to balancing. It therefore covers
the whole line; a bounded cycle cannot be embedded into this line as a
pure polyhedral graph. A component supplied already as a full line has
the same conclusion. Its tangent lattice is the saturated primitive line
lattice \(\mathbb{Z}v\), independently of the ambient lattice's
complementary directions.

For the all-coface PCC tests on a cylinder \(G\times\mathbb{R}\), any
convex cell projects to a convex subset of the graph. A segment cannot
pass a branched vertex since it would need an opposite pair there. Thus
its graph projection is in a single straight chain through bivalent
subdivisions or in a vertex. If the centre's graph coordinate is in a
straight-chain interior, all cofaces containing the common face remain
in that chain and the graph test is ordinary one-variable convexity. At
a branched centre \(v\), each coface projects into one incident straight
chain; every endpoint displacement has the form \(t_i v_e\),
\(t_i\geq0\), up to the next branched vertex. Grouping the barycentric
displacements gives \[
 b_e=\sum_{i\text{ in }e}\alpha_i t_i,
 \qquad\sum_e b_ev_e=0,\qquad b_e=c r_e,\ c\geq0.           \tag{9}
\] This argument uses actual cofaces on arbitrary rational structures,
not merely a prescribed product subdivision. Distant endpoints and
inserted bivalent vertices are included. Zero graph displacements cause
no contribution.

\subsection{Auxiliary rational-reference localization: whole-space
growth}\label{auxiliary-rational-reference-localization-whole-space-growth}

In this auxiliary construction only, write \(\eta\) for any
rational-coefficient positive-growth actual-PC reference. Its elementary
construction will be applied separately to every component \(\gamma_j\)
of (FC6); it is not a replacement of the prescribed real reference.

The following whole-space growth estimate applies to every product, with
no rank condition. Choose any rational centre \(A,B\); in the rank-one
two-branch construction these will be original branched vertices. All
objects in this construction remain on the \textbf{entire original
product}. Choose ambient lattice bases and put \[
\rho(p,q)=\|p-A\|_1+\|q-B\|_1,\qquad
\rho_0(r)=\|r_1\|_1+\|r_2\|_1.                               \tag{28}
\] The second is the \textbf{homogeneous recession norm}, not the
translated function \(\rho\). Refine the original product and \(\eta\)
into a common finite rational structure \(\Sigma\), splitting any
full-line carrier pieces if needed. Each nonempty cell is pointed.
Indeed a lineality vector \(r\ne0\) and \(-r\) would both be recession
vectors of an affine reference cell, and its affine slopes on them are
opposite; they cannot both be strictly positive. Any subcell inherits
positive slopes on its nonzero recession vectors.

By (8), write every cell as the convex hull of finitely many rational
vertices \(v_i\) plus the cone of finitely many nonzero rational
recession generators \(r_j\). Its affine slope
\(\eta_\sigma^{\mathrm{lin}}(r_j)\) is a positive rational number.
Select one rational \(m>0\) no larger than any of the finitely many
ratios \[
       \eta_\sigma^{\mathrm{lin}}(r_j)/\rho_0(r_j).            \tag{29}
\] If there are no recession generators, choose any rational \(m>0\).
For \(x=\sum_i\beta_iv_i+\sum_jt_jr_j\), \(\beta_i,t_j\geq0\),
\(\sum_i\beta_i=1\), affinity and the norm triangle inequality give \[
\begin{split}
\eta(x)&\geq\min_i\eta(v_i)+m\sum_jt_j\rho_0(r_j),\\
\rho(x)&\leq\max_i\rho(v_i)+\sum_jt_j\rho_0(r_j).
\end{split}                                                        \tag{30}
\] Choose rational \[
B=\max\bigl(0,\ \max_{\sigma\in\Sigma}
                 [m\max_{v\in\operatorname{Vert}\sigma}\rho(v)
                        -\min_{v\in\operatorname{Vert}\sigma}\eta(v)]\bigr).
\] The finite vertex sets are nonempty, and all values in this
expression are rational. Equations (30) prove on \textbf{all}
components, bounded cells and ends that \[
                         \eta\geq m\rho-B.                 \tag{31}
\] This derives the needed coercivity from the exact recession
hypothesis, rather than assuming a global growth theorem or using a
conveniently chosen reference.

\section{Complete multivariate cylinder
reconstruction}\label{complete-multivariate-cylinder-reconstruction}

In this unit graph weights are the positive rational normalized weights.
The main text proves the componentwise real-weight normalization and
restores the entire operator scalar. The divisor referred to as (5) here
is the actual formula (D) proved in every dimension above. All
primitive-normal, locality, balancing, refinement and commutativity uses
in this unit refer to that reconstructed algebra. This full multivariate
proof is mathematical evidence assessed here, not a theorem grant from
its predecessor.

For the inherited labels in this reconstruction, (1) denotes the sole
positive rank-one star relation, (2) denotes contact identity (C),
(3) denotes raw-envelope definition (E), and (5) denotes the full
divisor formula (D). These are aliases for the same formulas above.

\subsection{Elementary polyhedral and convex
facts}\label{elementary-polyhedral-and-convex-facts}

We record proofs of the finite-dimensional facts used subsequently. A
cone defined by finitely many rational inequalities has finitely many
rational generators. Remove its rational lineality first. In the pointed
quotient the sum of the defining nonnegative forms is positive away from
zero; its section of value one is bounded, since an unbounded sequence
normalized on the sphere would give a nonzero vector of value zero. A
bounded polyhedron is the convex hull of its finitely many vertices:
within the minimal face of any nonvertex point, a line through the point
ends on smaller-dimensional faces, and induction reduces to vertices.
Vertices solve full-rank subsystems, hence are rational for rational
data. Lift back and add both signs of a rational lineality basis.

Homogenization applies this assertion to a pointed rational polyhedron
and gives a nonempty finite point set and a finite ray set with \[
 \sigma=\operatorname{conv}(A_\sigma)+\operatorname{cone}(R_\sigma),
                    \qquad A_\sigma,R_\sigma\text{ rational}.       \tag{4}
\] The same assertions, without rationality of coordinates, hold for
real inequality data. Eliminating one variable from nonstrict finite
inequalities compares every lower bound with every upper bound,
retaining conditions independent of the variable; a one-sided family
imposes no further condition. Induction proves that projections of
polyhedra are polyhedra. Coefficients of variables remain rational if
they start rational, even if the constant terms are real. In particular
finitely generated convex-hull-plus-cone sets are closed. These
arguments prove the finite-polyhedron and elimination facts used here
rather than invoking an external polyhedral theorem.

All separations used below follow from nearest points in
finite-dimensional Euclidean space. A point outside a nonempty closed
convex set has a nearest point by restricting to a sufficiently large
compact ball. Differentiating squared distance along a segment in the
set gives a separating normal. For a disjoint compact convex set and
closed convex set, distance is positive and attained by the same compact
restriction; differentiation in both sets strictly separates them by
their nearest-point difference. For a cone the finite supremum of the
separating functional is zero. This proves the point/cone and
compact-polytope/cone separations used below, without importing a
separation theorem.

A finite rationally convex function on \(\mathbb{Q}^n\) has a unique
continuous convex completion to \(\mathbb{R}^n\). Corner values bound it
above on rational boxes. Reflection about a rational centre gives a
lower bound on a smaller box: \(F(y)\geq2F(c)-F(2c-y)\). Choose bounds
\(L\leq F\leq M\) on a rational box \(c+[-2R,2R]^n\). For rational
\(y,z\in c+[-R/2,R/2]^n\), set \(\delta=\|z-y\|_\infty\). If
\(\delta>0\), the rational factor \(T=R/\delta\geq1\) puts \(y+T(z-y)\)
in the larger box. Jensen gives \(F(z)-F(y)\leq(M-L)\delta/R\);
reversing the pair gives the opposite bound. This proves local Lipschitz
completion; density proves convexity and uniqueness. A finite convex
function on all of \(\mathbb{R}^n\) is continuous by the same box
argument.

It has an affine support at each point. Separate a point just below its
closed epigraph by the normal at a nearest epigraph point; nearest
points exist by compact restriction, and the separating inequality
follows by differentiating squared distance along line segments. The
vertical normal coefficient is positive, since a nonzero horizontal
normal cannot bound all of \(\mathbb{R}^n\). Normalize it. Evaluating
the resulting affine minorants at \(x\pm Re_i\) bounds their slopes as
the below-epigraph point approaches \((x,F(x))\). A convergent
subsequence yields a support touching at \(x\). This also covers
\(n=0\).

For a compact polyhedron \(D\), its outward normal cone at \(p\) is the
cone of the active defining inequality normals. Indeed membership in
this cone implies the normal inequality. If a normal were outside that
closed finitely generated cone, separation would provide a direction
satisfying the active inequalities and increasing the proposed normal; a
sufficiently small step satisfies all inactive inequalities too, a
contradiction. This proves both the assertion and its rational-generator
version when inequality normals are rational, irrespective of real
constants.

Finally, if \(a(p)=\min_i(c_i+g_i(p))\) is concave polyhedral, its
supergradient set at \(p\) is the convex hull of the active \(g_i\). The
directional derivative is their minimum. Separation proves the
supergradient assertion: an alleged supergradient outside that convex
hull gives a violating direction. For a finite sum of such functions
with nonnegative coefficients, the supergradient set is the
corresponding Minkowski sum, directly by expanding the finite minima.
These facts hold also at zero coefficients.

\subsection{Scalar graph facts proved without PC
locality}\label{scalar-graph-facts-proved-without-pc-locality}

The strictly positive vector in (1) implies that the full kernel of the
direction map is one-dimensional. Otherwise a small perturbation of that
vector by another kernel vector stays positive and contradicts (1). At a
bivalent vertex the actual primitive directions are opposite and the
weights equal. At any vertex of larger valence there is no opposite
pair: such a pair would give a second nonnegative relation. Subdivision
at rational edge points therefore preserves (1).

A convex polyhedron contained in \(G\times E\) projects to a convex
subset of \(G\). A straight segment cannot pass a branching vertex,
since there is no opposite pair there. Thus the graph projection is in a
single straight edge chain through bivalent vertices, or is a vertex.
This applies to every permitted rational structure, not just our chosen
subdivision.

\textbf{Scalar lemma.} Suppose \(k\) is finite on rational graph points
and satisfies rational graph PCC inequalities with rational
coefficients. It has a continuous completion, is convex on every edge,
has finite outgoing derivatives \(s_{v,e}\), and satisfies
\(\sum_e w_es_{v,e}\geq0\). Conversely a continuous function convex on
the edges with these inequalities, when added to any affine function of
\(E\), satisfies every actual cylinder PCC, with real coefficients and
all real centres.

For the first assertion, complete the rational Jensen function on each
edge interior. A one-sided limit at an endpoint is finite and at most
the specified endpoint value: an endpoint secant bounds it above and an
interior secant bounds it below. For rational small \(t>0\), the points
\(v+t v_{v,e}\) with coefficients \(w_e/\sum w_e\) are an actual vertex
PCC. Its inequality and the endpoint upper bounds force all limits to
equal \(k(v)\). Every outgoing derivative is bounded above by a finite
secant. Dividing the same PCC inequality by \(t\) prevents any
derivative being \(-\infty\), and gives the weighted derivative
inequality.

For the converse, edgewise convexity and the bivalent inequalities give
ordinary convexity on every straight chain. At a branching centre \(v\),
group the graph displacements of any cylinder PCC as \(t_i v_{v,e}\),
with \(t_i\geq0\), using the preceding structural observation. Their
weighted sum is zero, so the group coefficients
\(b_e=\sum_{i\text{ in }e}\alpha_it_i\) equal \(cw_e\), \(c\geq0\).
Convexity gives \(k(q_i)-k(v)\geq t_i s_{v,e}\). Summing gives a
nonnegative result. At a centre in a straight chain the assertion is
ordinary convexity. Terms with zero graph displacement and the affine
\(E\) terms are handled by the barycentre identity. This proves every
actual PCC directly, including remote endpoints, irrational centres, and
structures crossing bivalent subdivision points. No open-cover locality
theorem is used.

For a raw cylinder competitor \(h\), its rational slice at a rational
graph point \(q\) has the finite convex completion just proved. For a
real dual covector \(p\in E^*\), put \[
 k_{h,p}(q)=\inf_{z\in E_{\mathbb{Q}}}(h(q,z)-p(z))\in[-\infty,\infty).             \tag{7}
\] It satisfies rational graph PCC inequalities: choose rational \(z_i\)
approaching the endpoint infima, apply the actual cylinder PCC at
\((q,\sum_i\alpha_i z_i)\), subtract \(p\), and take the infimum.
Rational coefficients ensure that this last affine coordinate is
rational. If one endpoint has infimum \(-\infty\), vary that endpoint
alone.

If (7) is finite at one rational graph point, it is finite at every
rational point in that connected component. From an edge-interior finite
centre, any other rational point on that edge can be an endpoint in a
rational combination by choosing the other endpoint on the opposite
side. At a finite vertex, any rational remote point \(v+t_ev_{v,e}\) on
a chosen incident edge belongs to a vertex PCC with positive rational
coefficients proportional to \(w_e/t_e\), choosing positive rational
distances within all other incident edges. No endpoint can have value
\(-\infty\). Propagate across the finite connected graph. This works
through cycles and all unbounded ends. The scalar lemma then applies to
(7).

\subsection{Finite partial-infimum data, with tails and recession
constraints}\label{finite-partial-infimum-data-with-tails-and-recession-constraints}

Work in one connected component, identify \(E\) rationally with its
lattice vector space, and take a common finite rational affine structure
for \(f=\gamma+u\) and \(\gamma\), refined by the original edge strips.
Each cell is pointed: a nonzero lineality direction would give two
opposite reference recession slopes, incompatible with positivity at
infinity. Formula (4) therefore applies.

For each cell ray with zero graph coordinate, write its affine obstacle
slope as \(s_r\) and its \(E\)-coordinate as \(r_E\), and define \[
 D=\{p\in E^*:p(r_E)\leq s_r\text{ for every such ray}\}.                    \tag{8}
\] All inequality normals are rational; constants may be real. On every
recession ray \(u\) has zero slope, by boundedness, so (8) is the same
for \(f\) and \(\gamma\). It contains 0 in its interior, because every
nonzero vertical ray has strictly positive reference slope. For \(m=0\),
\(D=\{0\}\) in the zero-dimensional dual space.

Set \(F_p(q)=\inf_{z\in E}(f(q,z)-p(z))\). For \(p\in D\) this is finite
at every \(q\) and its infimum is attained at a vertex of some cell
fibre. Indeed, the fibre is pointed and its recession cone is generated
by the zero-graph cell rays: graph recession is zero on bounded strips
and nonnegative on unbounded strips, so a zero-graph combination cannot
use any positive-graph ray. The objective is nonnegative on this
recession cone, and (4) for the fibre gives a minimizing vertex. There
are finitely many cells. If \(p\notin D\), a violating rational ray
based at a rational point gives \(F_p(q)=-\infty\) at a rational graph
point. Thus (8) is precisely the common finite domain required here.

The set \(D\) is compact. At any fixed rational graph vertex, the
finite-PA function \(f(q,z)\) is bounded above by \(M\|z\|+B\), for some
finite \(M,B\) (use finitely many affine pieces and the dual norm). If
\(\|p\|>M\), a direction with \(p(z)>M\|z\|\) makes the infimum
\(-\infty\). Finiteness for \(p\in D\) bounds \(D\). This proof does not
require \(f\) convex.

Subdivide each bounded edge at the graph coordinates of all its spatial
cell vertices, and each end at those coordinates and its original base.
This is finite and rational. Between consecutive coordinates a fibre
vertex lies on a spatial one-face: the minimal face containing it loses
one dimension under the slicing hyperplane, so has dimension at most
one. A one-face traversing the interval has rational form
\(z(t)=z_i+t b_i\), and obstacle height real-affine in \(t\). Its
feasibility across that open interval is constant, because changes occur
only at projected spatial vertices. Therefore on each such interval \[
 F_p(t)=\min_i\{A_i(p)+t B_i(p)\},                      \tag{9}
\] where the finite index set is independent of \(p\); \(A_i,B_i\) are
affine in \(p\), with rational linear coefficients and real constants.
At a node \(v\), slicing the cells directly gives \(F_p(v)\) as a finite
minimum of affine functions of \(p\), again with rational gradients and
real constants.

For clarity, (9) extends concavely and continuously to each finite
interval endpoint. Its candidate fibre vertices have finite limits
there, so their limiting objective values are at least \(F_p(v)\).
Conversely, fixing any finite \(z\) and using continuity of \(f(t,z)\)
gives \(\limsup_{t\to v}F_p(t)\leq f(v,z)-p(z)\); then take its infimum
over \(z\). This proves equality of the limits. It also proves
branchwise continuity at a graph vertex. The possibility of a different
endpoint fibre type is not discarded.

After the last projected vertex on an end, translate its primitive
parameter to start at 0. Formula (9) holds for all \(t\geq0\), and \[
 \theta_e(p):=\lim_{t\to\infty}F_p(t)/t=\min_i B_i(p),\qquad
 F_p(t)-\theta_e(p)t\text{ is nondecreasing}.             \tag{10}
\] Every \(B_i-\theta_e\geq0\) proves the monotonicity term by term. In
particular the line from any value at most \(F_p(0)\) with slope
\(\theta_e(p)\) lies below the entire tail obstacle. All recession
constraints and finite endpoints have now been retained.

\subsection{The graph incidence polyhedron and its greatest
assignment}\label{the-graph-incidence-polyhedron-and-its-greatest-assignment}

Let \(V\) be the finite node set just obtained. For fixed \(p\in D\),
introduce node values \(a_v\), and impose \[
 a_v\leq F_p(v),\qquad
 L_v(a,p):=\sum_{e=vw\text{ bounded}}\frac{w_e}{\ell_e}(a_w-a_v)
                +\sum_{e\text{ tail at }v}w_e\theta_e(p)\geq0.             \tag{11}
\] Here \(\ell_e>0\) is the actual rational primitive length. All
bounded edges are present in this one system, including each edge of
every cycle. Do not integrate a graph direction into a larger lift or
unwind the graph.

There is a feasible assignment. The raw PC function \(\gamma-C\),
\(C=\|u\|_\infty\), lies below \(f\). Its partial infimum is finite for
\(p\in D\), by (8) applied to \(\gamma\), and the scalar lemma applies.
Replace its restriction on each bounded interval by its endpoint chord,
and on each tail by the line of slope \(\theta_e(p)\) from its endpoint
value. Convexity and the upper bound \(F_p\) show that the chords lie
below the concave data (9); the tail bound follows from (10). The chord
outgoing slopes at both endpoints are at least the old outgoing
derivatives. On a tail the asymptotic upper bound by (9) forces every
secant slope of the old convex function to be at most \(\theta_e(p)\),
so the new outgoing slope is also at least the old one. Thus all
inequalities in (11) hold.

The componentwise maximum of two feasible assignments is feasible. At a
node attaining its value from one assignment, each neighboring value has
only increased, so \(L_v\) has not decreased; the caps also persist.
Choose a feasible witness \(b\). Maximizing \(\sum_v a_v\) after taking
componentwise maximum with \(b\) restricts to the compact box
\(b_v\leq a_v\leq F_p(v)\), and hence attains a maximum. A feasible
assignment exceeding that maximum in any coordinate would increase this
objective by taking componentwise maximum. Consequently there is a
unique componentwise greatest assignment, denoted \(a_v(p)\).

It also dominates every scalar function \(k\) obtained from (7) that
lies below \(F_p\). The chord-and-tail replacement just used dominates
\(k\) everywhere by convexity and the tail slope bound, and its node
vector is feasible for (11). Therefore it is at most the interpolation
of the greatest assignment. Define that interpolation as \(k_p(q)\):
endpoint-linear on every bounded interval and \(a_v(p)+t\theta_e(p)\) on
every tail. It is continuous, scalar PC by the scalar lemma, and
\(k_p\leq F_p\) everywhere.

Every \(a_v(p)\) is a finite concave piecewise-affine function on \(D\),
with rational gradients and real constants. To see this without an
unproved parametric-LP assertion, expand all caps \(F_p(v)\) into their
finite affine minima. Expand each tail minimum \(\theta_e\) in the last
inequalities of (11): the weighted sum of minima is the minimum over all
choices of tail pieces. Thus (11) is one finite linear system in
\((a,p)\) with rational variable coefficients and real constants,
together with (8). Eliminate all \(a_w\), \(w\ne v\), by the elementary
elimination procedure. For fixed \(p\in D\), its section in \(a_v\) is
exactly \((-\infty,a_v(p)]\), because simultaneously lowering all node
values preserves the Laplacians and caps. Projected inequalities
involving \(a_v\) must therefore be upper bounds; zero-coefficient
inequalities concern \(D\). There is at least one upper bound from the
original cap. Dividing positive rational coefficients gives \[
 a_v(p)=\min_j(c_{v,j}+g_{v,j}(p)),\qquad g_{v,j}\text{ rational}.             \tag{12}
\] This also proves continuity on all of \(D\), including its boundary.
Tail slopes already have the same property. No assertion that the dual
vertices are rational is made.

\subsection{Actual raw envelope and finite rational spatial
domains}\label{actual-raw-envelope-and-finite-rational-spatial-domains}

Define \[
 Q(q,z)=\sup_{p\in D}\{k_p(q)+p(z)\}.                                    \tag{13}
\] It is finite: \(D\) is compact and all node values and tail slopes
are continuous. For each fixed \(p\), \(k_p(q)+p(z)\) satisfies every
actual PCC by the scalar lemma, so its supremum does too
\(take the supremum of each Jensen inequality\). Since \(k_p\leq F_p\),
\(Q\leq f\).

For every raw competitor \(h\leq f\), and every rational \((q,z)\), take
a touching affine support of the convex completion of the rational slice
\(h(q,\cdot)\). Its covector \(p\) makes (7) finite and
\(h(q,z)=k_{h,p}(q)+p(z)\). Propagation proves finiteness throughout the
component; a violating ray would contradict \(h\leq f\), so \(p\in D\).
The preceding greatest-assignment argument gives \(k_{h,p}\leq k_p\).
Hence \(h(q,z)\leq Q(q,z)\) at all rational carrier points.

The function \(Q\) is continuous. Partition \(D\) into the common finite
polyhedral chambers of all affine minima in (12) and (10). Each chamber
is compact. On it \(k_p(q)+p(z)\) is affine in \(p\) on a fixed graph
interval, so its maximum is attained at a chamber vertex. Let \(J\) be
the finite union of all these vertices, with no rationality assumption
on their coordinates. Formula (13) on each edge strip is the maximum of
the corresponding finitely many real-affine functions of \((q,z)\). At
graph nodes the formulas agree. This proves continuity on all of \(X\).

To extend the raw upper bound to a real point \(x\), choose any finite
rational structure and its minimal cell containing \(x\). In a
positive-dimensional cell, shrinking simplices with rational vertices in
its relative interior enclose \(x\); they can be obtained by rationally
perturbing the vertices of a small enclosing simplex in the rational
affine hull. Their real barycentric coefficients give actual same-cell
PCCs. Raw Jensen and continuity of \(Q\) imply \(h(x)\leq Q(x)\). A
zero-dimensional cell point is already rational. No raw-competitor
continuity has been used.

Apply this to \(h=\gamma-C\). Together with \(Q\leq f\leq\gamma+C\), it
gives the two-sided growth bound. Consequently \(Q\) itself belongs to
the competitor set (3), and all its competitors are at most \(Q\). Thus
\[
                        P_\gamma(u)=Q.                                  \tag{14}
\] Both the global lower competitor and global upper comparison are
exact, and no local raising or unproved growth gluing is invoked.

Real-affine maxima alone would not prove rational spatial domains. Here
there is an additional rational-gradient mechanism. Fix \(p_0\in J\).
Let \(A_v(p_0)\) be the supergradient polytope of \(a_v\) there, and let
\(B_e(p_0)\) be that of \(\theta_e\). Their vertices are rational by
(12). Let \(N_D(p_0)\) be the outward normal cone of \(D\), generated by
rational normals by (8).

For any finite concave polyhedral \(b\), the condition that \(p_0\)
maximize \(b(p)+p(z)\) over \(D\) is \[
                    z\in-\partial^+b(p_0)+N_D(p_0).                      \tag{15}
\] For sufficiency, a supergradient \(g\) with \(g+z\in N_D\) gives the
global maximum inequality. For necessity, if \((z+\partial^+b)\cap N_D\)
is empty, separate the compact polytope from the closed polyhedral cone.
The separating direction \(r\) has \((g+z)(r)>0\) for every
supergradient and has \(n(r)\leq0\) for every normal. It is a feasible
tangent direction to \(D\); a sufficiently small step is feasible. The
directional derivative, the minimum of these positive quantities,
increases the objective, a contradiction. This supplies the exact
optimality argument, including boundary faces.

On a bounded interval \(t\in[0,\ell]\), \(b=(1-t/\ell)a_v+(t/\ell)a_w\),
so (15) describes the full maximization domain of the affine plane
indexed by \(p_0\) as \[
 z\in-(1-t/\ell)A_v(p_0)-(t/\ell)A_w(p_0)+N_D(p_0).                       \tag{16}
\] It is the convex hull of the rational endpoint point sets
\((0,-\operatorname{Vert}A_v)\), \((\ell,-\operatorname{Vert}A_w)\),
plus the cone of \((0,n)\) for rational generators of \(N_D\). On a tail
its domain is \[
 t\geq0,\qquad z\in-A_v(p_0)-tB_e(p_0)+N_D(p_0),                         \tag{17}
\] generated by the rational points \((0,-\operatorname{Vert}A_v)\) and
rational rays \((1,-\operatorname{Vert}B_e)\), \((0,n)\). These
descriptions remain exact at \(t=0,\ell\), at zero coefficients, and at
degenerate chambers. They prove that every maximization domain is a
rational spatial polyhedron, though its affine value has real
coefficients.

The maximization domains of a finite affine maximum, including all their
faces and removing duplicates, form a complex: intersections are
equality faces of the respective maxima regions. Its faces are rational
because these regions are rational polyhedra. At a graph node (16)--(17)
have exactly the same fibre domain for each \(p_0\), independent of the
neighboring interval, so the strip complexes glue. This gives a finite
rational spatial complex for \(P\) on the actual carrier. Moreover (16)
has vertices only at its two graph endpoints, and (17) only at its base;
faces introduce no new vertices. Thus every possible zero-dimensional
domain face of this complex lies over a node of \(V\). Edge interiors
cannot carry a recursive zero-cycle coefficient. This conclusion
includes all bounded-edge interiors and all tails, rather than omitting
their lower-face terms.

\subsection{PAPC against all local real cycles and actual
regularization}\label{papc-against-all-local-real-cycles-and-actual-regularization}

Let \(z\geq0\) be any local real balanced cycle of any positive
dimension on an ordinary open \(U\subset X\), and refine to make \(P\)
affine. At a tested facet \(\tau\), choose a rational
\(x\in U\cap\operatorname{relint}\tau\). The cone of possible
nonnegative incident weights balanced modulo \(N_\tau\) is a rational
polyhedral cone, so is generated by rational vectors, by the elementary
cone proof above. The divisor coefficient is linear in those weights. It
suffices to test a nonnegative rational incident vector \(c_\sigma\).

Write \(S=\sum c_\sigma>0\),
\(B=\sum c_\sigma n_{\sigma/\tau}\in N_{\tau,\mathbb{Q}}\), and use the
rational inward vectors \(n_{\sigma/\tau}-B/S\). They need not be
integral normals; they have the same positive quotient directions. For
sufficiently small rational \(\varepsilon>0\), the rational points
\(x_\sigma=x+\varepsilon(n_{\sigma/\tau}-B/S)\) belong to
\(\sigma\cap U\). With coefficients \(c_\sigma/S\) their barycentre is
\(x\), giving an actual PCC. Its Jensen inequality is \[
 0\leq\sum c_\sigma[P(x_\sigma)-P(x)]
       =\varepsilon\left(\sum c_\sigma P_\sigma(n_{\sigma/\tau})-P_\tau(B)\right).
\] This is exactly (5). The zero-weight case is zero.
Rational-generation then proves the inequality for every nonnegative
real incident vector, and hence for every real local cycle. Thus \(P\)
is PAPC in the full required sense. Balancing proved above permits
recursion, and every successive divisor is nonnegative. Its full top
cycle is a finite nonnegative atomic measure on the rational domain
vertices. The sequence \(P+1/j\) is PAPC in the same two-sided growth
class and decreases uniformly to \(P\); its divisors agree with those of
\(P\). This establishes the actual operator domain and regularization
without invoking source regularity assertions.

\subsection{Two complete recursive vanishing
lemmas}\label{two-complete-recursive-vanishing-lemmas}

\textbf{Quotient vanishing.} On a local \(d\)-cycle, suppose that after
subtracting an ambient affine function the PA germ is pulled back by a
rational linear map to a function on a complete rational fan in a space
of dimension \(q<d\). Its full top divisor is zero. Let \(T_j\) be the
union of quotient fan faces of codimension at least \(j\). The first
divisor vanishes off the preimage of \(T_1\), since the function is
ambient affine there. Inductively, near the relative interior of a
quotient face of codimension \(j-1\), the preceding cycle is supported
in that face's preimage, where the function is affine on the entire
supporting subspace. Extend this affine function to the ambient space;
balancing makes its next divisor zero. Thus the \(j\)-th divisor is
supported in the preimage of \(T_j\). Since \(T_{q+1}\) is empty, the
top divisor vanishes. Rational common refinements and the local formula
(5) justify this argument on every lower face, without transversality or
a top-face-only measure formula.

\textbf{Lineality vanishing.} If the local carrier fan and residual
function-domain fan have common nonzero rational lineality, choose their
common generalized fan retaining that lineality. Every face contains it;
there is no zero-dimensional face. Every recursive divisor lives on that
fan, so the top zero-cycle is zero. Refinement cannot change this by the
proved refinement identity. Ambient affine terms can be removed first.

\subsection{Full contact at cyclic junction
fibres}\label{full-contact-at-cyclic-junction-fibres}

Let \(x=(v,z_0)\) be any possible atom of the finite rational spatial
complex for \(P\). By (16)--(17), \(v\in V\) and \(z_0\) is rational.
Suppose \(P(x)<f(x)\). Define its compact exposed dual face \[
 C_x=\{p\in D:a_v(p)+p(z_0)=P(v,z_0)\}.                                  \tag{18}
\] It is nonempty and convex. Indeed \(a_v\) is concave, and the
displayed quantity never exceeds its maximum. Formula (12) describes
(18) by inequalities with rational normals \(g_{v,j}+z_0\) and real
constants, together with (8). Its affine-hull direction is therefore
rational: the affine hull is cut out by the inequalities that are
equalities throughout the polyhedron, as seen by choosing a
relative-interior point and varying in its remaining feasible
directions. Their homogeneous coefficients are rational. On \(C_x\),
\(a_v(p)=P(x)-p(z_0)\) is affine.

For each incident bounded edge \(e=vw\), directed away from \(v\), and
each tail, put \[
 s_e(p)=\frac{a_w(p)-a_v(p)}{\ell_e}\quad\text{or}\quad s_e(p)=\theta_e(p).
                                                                         \tag{19}
\] On \(C_x\) these are concave continuous polyhedral functions: \(a_v\)
is affine there. The vertex inequality says \[
                           \sum_e w_es_e(p)\geq0.                         \tag{20}
\] The local germ of \(P-P(x)\) in the outgoing strip with coordinates
\(t\geq0, \zeta=z-z_0\) is exactly \[
                  h_e(t,\zeta)=\max_{p\in C_x}\{p(\zeta)+t s_e(p)\}.       \tag{21}
\] To justify this, use the fixed finite set \(J\) of common dual
chamber vertices. All inactive planes at \(x\) have strictly positive
gap, and therefore stay inactive in a sufficiently small neighbourhood.
Active planes have their dual points in \(C_x\). Conversely every point
of \(C_x\) belongs to a chamber and is a convex combination of its
vertices; equality in the exposed maximum forces those vertices to be
active, and on that chamber all terms in (19) are affine. Thus (21)
follows exactly. This is a local finite-germ identity, not an assumption
of raw continuity or joint-convex reconstruction.

If \(\dim C_x<m\), let \(T\subset E\) annihilate the rational direction
of \(C_x\). It is nonzero and rational. Subtract the ambient affine
\(p_0(z)\), with \(p_0\in C_x\). Formula (21) is invariant under \(T\)
in every strip; the carrier star has \(E\), hence \(T\), as lineality. A
rational function-domain fan retaining \(T\) exists: each coarsest
linearity region is a convex union of finitely many rational tangent
cones of the already constructed spatial complex, and hence is a
rational cone; its equations are invariant under \(T\). Quotient these
cones by \(T\) and lift their common rational refinement. The lineality
lemma proves that the full coefficient at \(x\) is zero. This case
includes all lower-dimensional exposed faces.

Now suppose \(\dim C_x=m\). For \(p\in\operatorname{int}C_x\), choose
\(\eta>0\) with \(p+\eta\overline B^*\subset C_x\). Formula (18) and
(13) give, for every \(y\in E\), \[
 P(v,y)-p(y)\geq P(v,z_0)-p(z_0)+\eta\|y-z_0\|.                           \tag{22}
\] Since \(f(v,\cdot)\geq P(v,\cdot)\) is continuous, \(f(v,y)-p(y)\) is
coercive and attains its infimum. If \(a_v(p)=F_p(v)\), that infimum
equals the first term on the right of (22). The strict inequality away
from \(z_0\) then forces the minimum at \(z_0\) and
\(f(v,z_0)=P(v,z_0)\), a contradiction. Therefore \[
                         a_v(p)<F_p(v)\quad(p\in\operatorname{int}C_x).    \tag{23}
\]

The greatest-assignment property now implies \(L_v(a(p),p)=0\). If this
Laplacian were positive, raise \textbf{only its scalar LP coordinate}
\(a_v\) by a sufficiently small \(\delta>0\). The cap persists by (23).
At \(v\) the Laplacian decreases by \(\delta\sum_{e=vw}w_e/\ell_e\), and
at every neighboring node it increases; tail slopes are fixed. Thus
feasibility persists, contradicting coordinatewise maximality. If there
are no bounded incident edges, raising the coordinate does not change
any Laplacian and gives the contradiction even without a positive
Laplacian; in that case (23) itself is impossible. This is a finite
scalar feasibility argument; it is not a local perturbation of \(P\)
requiring global PC gluing.

Continuity extends the equality to all of \(C_x\): \[
                              \sum_e w_es_e=0\quad\text{on }C_x.           \tag{24}
\] Since each \(s_e\) is concave and all weights are strictly positive,
every \(s_e\) is affine on \(C_x\). Indeed the nonnegative Jensen
deficits of the individual concave functions have weighted sum zero and
must each vanish. On the full-dimensional \(C_x\), choose an interior
point off the finitely many affine-piece walls. Formula (12) and the
rational lengths show that these affine expressions have rational
\(p\)-gradients. Write \[
                    s_e(p)=\alpha_e+p(\beta_e),\qquad\beta_e\in E_{\mathbb{Q}}.
                                                                         \tag{25}
\] The constants \(\alpha_e\) may be real. Equation (24) says
\(\sum w_e\alpha_e=0\) and \(\sum w_e\beta_e=0\).

Let \(S_v\) be the actual rational span of the primitive star
directions. Its only relation is (1). Consequently there is a real
linear form \(A:S_v\to\mathbb{R}\) with \(A(v_{v,e})=\alpha_e\), and a
rational linear map \(B:S_v\to E\) with \(B(v_{v,e})=\beta_e\).
Rationality follows by solving the rational direction matrix for the
rational values, not by assuming ray-generated lattice saturation.
Formula (21) is therefore the restriction to the entire star of \[
                       A(s)+\operatorname{supp}_{C_x}(\zeta+B(s)).        \tag{26}
\] After removing \(A+p_0\circ(\zeta+B(s))\), it factors through the
rational map \((s,\zeta)\mapsto\zeta+B(s)\) to the \(m\)-dimensional
complete space \(E\). The support function of \(C_x-p_0\) has a complete
rational normal fan, because (18) has rational inequality normals; its
normal cones are generated by active normals as proved above. Apply
quotient vanishing with \(q=m<d=m+1\). The entire recursive coefficient
at \(x\) is zero, with actual primitive indices and every intermediate
lower-face contribution retained.

The proof covers every possible atom: original cyclic junctions,
inserted bivalent nodes, all fibre domain faces, and tail bases.
Non-node edge interiors have already been excluded by the
rational-domain generator argument. No tree branch decomposition, cycle
cutting, globally affine lift, or independent support extension is used.
For \(m=0\) the same argument reduces to the known graph contact
mechanism; the new assertion is every \(m>0\) with arbitrary
nonseparable data and cycles. Since the measure is finite nonnegative
and every noncontact coefficient is zero, (2) follows. The integrand is
bounded by the two-sided growth inequalities, and no boundary atom is
introduced: the PA operator is the finite zero-cycle on \(X\), and its
explicit constant-shift regularization has exactly the same measure.
Infinity and all ends have been handled in (8)--(11).

For finitely many connected components, construct \(D\), the LP, and
\(P\) separately. A PCC and a local cycle stay in the component of their
incident cell; the finitely many measures add and a uniform growth
constant is available. This finishes the stated graph-cylinder theorem.

\section{Detailed original-lattice bad-pair mass
computation}\label{detailed-original-lattice-bad-pair-mass-computation}

This computation is reconstructed at the original singular product star.
Its positive scale called epsilon in the final two paragraphs is kappa
in the main theorem. The all-stage tensor formula (A3), supplied above,
is an additional proof and is not inferred merely from this two-stage
computation.

\subsection{Complete primitive-normal junction
computation}\label{complete-primitive-normal-junction-computation}

In the original product quadrant \((s v_i,t w_j)\), \(s,t\geq0\), put
\(q_i=1/S_k\) if \(i\in I_k\), \(p_j=1/T_l\) if \(j\in J_l\). Directly
from the three specified tripod directions, \[
 H=c_kq_i s+c_lp_j t\quad(k\ne l),\qquad
 H=c_k|q_i s-p_jt|\quad(k=l).                              \tag{40}
\] Let \(a_i=(v_i,0)\), \(b_j=(0,w_j)\). The actual saturated plane
lattice is \(\mathbb{Z}a_i\oplus\mathbb{Z}b_j\): a vector in this plane
belongs to the direct product ambient lattice exactly when its two
primitive factor-line coefficients are integral. Factor star spans may
be nonsaturated in their ray generators; this does not alter these
actual plane lattices.

Off-group quadrants are affine and have index-one axis cones. For a
same-group pair take a positive integer denominator \(D\), \(Q=Dq_i\),
\(P=Dp_j\), \(g=\gcd(Q,P)\). Its zero kink ray has actual primitive
vector \[
                            d_{ij}=(P/g)a_i+(Q/g)b_j.        \tag{41}
\] The integer coefficients are coprime; primitiveness follows also in
the whole ambient lattice because both factor vectors are primitive. The
two actual saturated cone indices are \[
                    \chi_{a_i,d_{ij}}=Q/g,\qquad
                    \chi_{b_j,d_{ij}}=P/g.                 \tag{42}
\] All top cone weights are \(\Omega_{ij}=ab\lambda_i\mu_j\), inherited
from the original graph product. At a kink, (6) and \(H(d_{ij})=0\) give
the full first weight \[
 D_H(d_{ij})=\Omega_{ij}\left(\frac gQ c_kq_i+\frac gP c_kp_j\right)
                             =2c_k\Omega_{ij}g/D.           \tag{43}
\] No denominator or index has been suppressed.

At an axis \(a_i\), the off-group neighbours are \(b_j\). For same-group
neighbours the representative quotient normal is
\(d_{ij}/(Q/g)=b_j+(p_j/q_i)a_i\). Therefore top balancing gives \[
 \sum\frac{\Omega_{ij}}{\chi}r_{ij}
       =ab\lambda_i\left(\sum_j\mu_jb_j+
                    \sum_{j\in J_k}\mu_j(p_j/q_i)a_i\right)
       =ab\lambda_i S_k a_i.                               \tag{44}
\] Since \(H(a_i)=c_k/S_k\), the off-group first term is
\(ab\lambda_i(4-c_k)\), and the tangential correction in (D), with
(44), subtracts \(ab\lambda_i c_k\). Thus \[
 D_H(a_i)=ab\lambda_i(4-2c_k),\qquad
 D_H(b_j)=ab\mu_j(4-2c_l).                                 \tag{45}
\] These axis rays and the kink rays (41) are all rays in the local
function fan and all are nonnegative. Its balancing can be checked
directly: kink contributions in the first projection for a fixed
\(i\in I_k\) sum to \(2c_kab\lambda_i v_i\), adding its axis
contribution gives \(4ab\lambda_i v_i\). The second projection similarly
gives \(4ab\mu_jw_j\). Both total sums vanish by the original relations.
This check retains every same-group pair even when multiple original
rays share a projected direction.

Apply the second divisor (7). Kink rays have zero value. The entire
first-axis contribution is \[
 ab\sum_k(4-2c_k)c_k\sum_{i\in I_k}\lambda_i/S_k=4ab,
\] and the second-axis contribution is \(4ab\). Hence the
\textbf{complete} coefficient at the chosen junction is \[
                          (H^2\cdot[X])\{(A,B)\}=8ab.              \tag{46}
\] This is the two-stage computation at the original singular product
junction, not a pushed-forward tripod calculation or a one apparent atom
claim.

By (J) the local germ of \(P-P(o)\) is \(\kappa H\). Constants
have zero divisor, and both stages of (5) are real linear in the
function, so the complete original-junction coefficient for \(P\) is \[
                    \mathrm{MA}_{\mathrm{poly},\mathbb{R}}(P)\{o\}
                                  =8\kappa^2ab.          \tag{47}
\] Other homogeneous cuts introduced by \(\gamma\) or the
carrier/refinement do not change this coefficient, by the proved
two-stage refinement compatibility. The calculation uses no pushforward
to a tripod, no missing-axis convention and no singular-face Alexandrov
substitution.

\section{Review, sources and reproducibility}

The independent Xean native Codex review used
\texttt{gpt-6.1-sol} with \texttt{xhigh} reasoning. Its full report was
read and bound to the exact immutable argument, manuscript and partial
status. The verdict covers this finite graph-product family. No
completion of the original arbitrary balanced-carrier problem is
claimed.

All new proofs and all three canonical native notes were read. Every
reused primary or comparison passage was checked for exact equality
with previously fully read text; all remaining lines were read directly.
All nineteen source snapshots were checked against their unchanged
origins, and all twelve selected raw source bodies were checked against
their exact inclusive LF coordinates and displayed copies.

The accompanying frozen argument retains the complete primary source
quotations and the existing one-dimensional comparison dossier. They
are comparison evidence. The complete mathematical dependency chain
is reconstructed in this manuscript; prior reviews are not theorem
premises. Integral source coefficients and the declared real extension
are distinguished throughout.

The four local recursive computations in dimensions three and five
were rerun from a fresh copy and matched the retained exact outputs.
Their full masses are \(64/49\), \(576000/26411\), \(64/567\), and
\(6400/6561\). They test specified local germs, all successive divisor
stages, all surviving orders and an exhausted repeated factor. They
do not enumerate global envelopes or prove universal contact integrals.
Those conclusions rest on the proofs above.

The packet contains the standalone source and PDF, complete frozen
argument, exact independent review, source provenance, and full finite
diagnostic program, input and output. Native editor verification and
local PDF export results are recorded separately. Same-model review,
typesetting and arithmetic are not formal verification or priority
certification.

\begin{thebibliography}{9}
\bibitem{BMPS}
A.~M.~Botero, E.~Mazzon and L.~Pille--Schneider,
\emph{Monge--Amp\`ere measures on balanced polyhedral spaces},
arXiv:2603.08664v1 (2026),
\url{https://arxiv.org/abs/2603.08664v1}.
\bibitem{AR}
L.~Allermann and J.~Rau,
\emph{First steps in tropical intersection theory},
arXiv:0709.3705v3 (2009),
\url{https://arxiv.org/abs/0709.3705v3};
\emph{Math.~Z.} \textbf{264} (2010), 633--670.
\end{thebibliography}

\end{document}
